%%%%%%%%%%%%%%%%%%%%%%%%%%%%%%%%%%%%%%%%%%%%%%%%%%%%%%%%%%%%%%%%%%%%%%%%%%%%
%DIF LATEXDIFF DIFFERENCE FILE
%DIF DEL before.tex   Tue Mar 10 00:07:21 2026
%DIF ADD after.tex    Tue Mar 10 00:07:21 2026
% AGUJournalTemplate.tex: this template file is for articles formatted with LaTeX
%
% This file includes commands and instructions
% given in the order necessary to produce a final output that will
% satisfy AGU requirements, including customized APA reference formatting.
%
% You may copy this file and give it your
% article name, and enter your text.
%
% guidelines and troubleshooting are here:

%% To submit your paper:
\documentclass{agujournal2019}

\usepackage{url} %this package should fix any errors with URLs in refs.
\usepackage{lineno}
\usepackage[inline]{trackchanges} %for better track changes. finalnew option will compile document with changes incorporated.
\usepackage{soul}
\usepackage{amsmath}
\usepackage{amssymb}


\linenumbers

\draftfalse

\journalname{JGR: Earth Surface, Confidential manuscript under review}
%DIF PREAMBLE EXTENSION ADDED BY LATEXDIFF
%DIF FONTSTRIKE PREAMBLE %DIF PREAMBLE
\RequirePackage[normalem]{ulem} %DIF PREAMBLE
\DeclareOldFontCommand{\sf}{\normalfont\sffamily}{\mathsf} %DIF PREAMBLE
\providecommand{\DIFadd}[1]{{\protect\color{blue} #1}} %DIF PREAMBLE
\providecommand{\DIFdel}[1]{{\protect\color{red}\sout{#1}}} %DIF PREAMBLE
%DIF COLOR PREAMBLE %DIF PREAMBLE
\RequirePackage{color} %DIF PREAMBLE
\providecommand{\DIFaddbegin}{\protect\color{blue}} %DIF PREAMBLE
\providecommand{\DIFaddend}{\protect\color{black}} %DIF PREAMBLE
\providecommand{\DIFdelbegin}{\protect\color{red}} %DIF PREAMBLE
\providecommand{\DIFdelend}{\protect\color{black}} %DIF PREAMBLE
\providecommand{\DIFmodbegin}{} %DIF PREAMBLE
\providecommand{\DIFmodend}{} %DIF PREAMBLE
%DIF FLOATSAFE PREAMBLE %DIF PREAMBLE
\providecommand{\DIFaddFL}[1]{\DIFadd{#1}} %DIF PREAMBLE
\providecommand{\DIFdelFL}[1]{\DIFdel{#1}} %DIF PREAMBLE
\providecommand{\DIFaddbeginFL}{} %DIF PREAMBLE
\providecommand{\DIFaddendFL}{} %DIF PREAMBLE
\providecommand{\DIFdelbeginFL}{} %DIF PREAMBLE
\providecommand{\DIFdelendFL}{} %DIF PREAMBLE
\newcommand{\DIFscaledelfig}{0.5}
%DIF HIGHLIGHTGRAPHICS PREAMBLE %DIF PREAMBLE
\RequirePackage{settobox} %DIF PREAMBLE
\RequirePackage{letltxmacro} %DIF PREAMBLE
\newsavebox{\DIFdelgraphicsbox} %DIF PREAMBLE
\newlength{\DIFdelgraphicswidth} %DIF PREAMBLE
\newlength{\DIFdelgraphicsheight} %DIF PREAMBLE
% store original definition of \includegraphics %DIF PREAMBLE
\LetLtxMacro{\DIFOincludegraphics}{\includegraphics} %DIF PREAMBLE
\newcommand{\DIFaddincludegraphics}[2][]{{\color{blue}\fbox{\DIFOincludegraphics[#1]{#2}}}} %DIF PREAMBLE
\newcommand{\DIFdelincludegraphics}[2][]{% %DIF PREAMBLE
\sbox{\DIFdelgraphicsbox}{\DIFOincludegraphics[#1]{#2}}% %DIF PREAMBLE
\settoboxwidth{\DIFdelgraphicswidth}{\DIFdelgraphicsbox} %DIF PREAMBLE
\settoboxtotalheight{\DIFdelgraphicsheight}{\DIFdelgraphicsbox} %DIF PREAMBLE
\scalebox{\DIFscaledelfig}{% %DIF PREAMBLE
\parbox[b]{\DIFdelgraphicswidth}{\usebox{\DIFdelgraphicsbox}\\[-\baselineskip] \rule{\DIFdelgraphicswidth}{0em}}\llap{\resizebox{\DIFdelgraphicswidth}{\DIFdelgraphicsheight}{% %DIF PREAMBLE
\setlength{\unitlength}{\DIFdelgraphicswidth}% %DIF PREAMBLE
\begin{picture}(1,1)% %DIF PREAMBLE
\thicklines\linethickness{2pt} %DIF PREAMBLE
{\color[rgb]{1,0,0}\put(0,0){\framebox(1,1){}}}% %DIF PREAMBLE
{\color[rgb]{1,0,0}\put(0,0){\line( 1,1){1}}}% %DIF PREAMBLE
{\color[rgb]{1,0,0}\put(0,1){\line(1,-1){1}}}% %DIF PREAMBLE
\end{picture}% %DIF PREAMBLE
}\hspace*{3pt}}} %DIF PREAMBLE
} %DIF PREAMBLE
\LetLtxMacro{\DIFOaddbegin}{\DIFaddbegin} %DIF PREAMBLE
\LetLtxMacro{\DIFOaddend}{\DIFaddend} %DIF PREAMBLE
\LetLtxMacro{\DIFOdelbegin}{\DIFdelbegin} %DIF PREAMBLE
\LetLtxMacro{\DIFOdelend}{\DIFdelend} %DIF PREAMBLE
\DeclareRobustCommand{\DIFaddbegin}{\DIFOaddbegin \let\includegraphics\DIFaddincludegraphics} %DIF PREAMBLE
\DeclareRobustCommand{\DIFaddend}{\DIFOaddend \let\includegraphics\DIFOincludegraphics} %DIF PREAMBLE
\DeclareRobustCommand{\DIFdelbegin}{\DIFOdelbegin \let\includegraphics\DIFdelincludegraphics} %DIF PREAMBLE
\DeclareRobustCommand{\DIFdelend}{\DIFOaddend \let\includegraphics\DIFOincludegraphics} %DIF PREAMBLE
\LetLtxMacro{\DIFOaddbeginFL}{\DIFaddbeginFL} %DIF PREAMBLE
\LetLtxMacro{\DIFOaddendFL}{\DIFaddendFL} %DIF PREAMBLE
\LetLtxMacro{\DIFOdelbeginFL}{\DIFdelbeginFL} %DIF PREAMBLE
\LetLtxMacro{\DIFOdelendFL}{\DIFdelendFL} %DIF PREAMBLE
\DeclareRobustCommand{\DIFaddbeginFL}{\DIFOaddbeginFL \let\includegraphics\DIFaddincludegraphics} %DIF PREAMBLE
\DeclareRobustCommand{\DIFaddendFL}{\DIFOaddendFL \let\includegraphics\DIFOincludegraphics} %DIF PREAMBLE
\DeclareRobustCommand{\DIFdelbeginFL}{\DIFOdelbeginFL \let\includegraphics\DIFdelincludegraphics} %DIF PREAMBLE
\DeclareRobustCommand{\DIFdelendFL}{\DIFOaddendFL \let\includegraphics\DIFOincludegraphics} %DIF PREAMBLE
%DIF COLORLISTINGS PREAMBLE %DIF PREAMBLE
\RequirePackage{listings} %DIF PREAMBLE
\RequirePackage{color} %DIF PREAMBLE
\lstdefinelanguage{DIFcode}{ %DIF PREAMBLE
%DIF DIFCODE_FONTSTRIKE %DIF PREAMBLE
  moredelim=[il][\scriptsize \sout]{\%DIF\ <\ }, %DIF PREAMBLE
  moredelim=[il][\sffamily]{\%DIF\ >\ } %DIF PREAMBLE
} %DIF PREAMBLE
\lstdefinestyle{DIFverbatimstyle}{ %DIF PREAMBLE
	language=DIFcode, %DIF PREAMBLE
	basicstyle=\ttfamily, %DIF PREAMBLE
	columns=fullflexible, %DIF PREAMBLE
	keepspaces=true %DIF PREAMBLE
} %DIF PREAMBLE
\lstnewenvironment{DIFverbatim}{\lstset{style=DIFverbatimstyle}}{} %DIF PREAMBLE
\lstnewenvironment{DIFverbatim*}{\lstset{style=DIFverbatimstyle,showspaces=true}}{} %DIF PREAMBLE
%DIF END PREAMBLE EXTENSION ADDED BY LATEXDIFF

\begin{document}

\title{Constraining the Distribution of 3D Fractal Structures in Mud Flocs}
\authors{
Brayden Noh\affil{1,2},
Justin A.~Nghiem\DIFdelbegin %DIFDELCMD < \affil{1}%%%
\DIFdelend \DIFaddbegin \affil{\DIFadd{1,3}}\DIFaddend ,
Kimberly Litwin Miller\affil{1},
Dongchen Wang\affil{1},
Michael P.~Lamb\affil{1}
}

\affiliation{1}{Division of Geological and Planetary Sciences, California Institute of Technology, Pasadena, CA 91125, USA}
\affiliation{2}{Department of Earth and Planetary Sciences, Harvard University, Cambridge, MA 02138, USA}
\DIFaddbegin \affiliation{3}{\DIFadd{St. Anthony Falls Laboratory, University of Minnesota, Minneapolis, MN 55414, USA}}
\DIFaddend 

\correspondingauthor{Brayden Noh}{braydennoh@fas.harvard.edu}


\begin{keypoints}
\item Aggregating floc data overestimates 3D fractal dimension due to structural diversity.
\item Stratifying by 2D \DIFdelbegin \DIFdel{fractal dimension reveals increasing $n_f$ with shear and decreasing $n_f$ with }\DIFdelend \DIFaddbegin \DIFadd{box-counting dimension reveals trends with shear velocity and }\DIFaddend organic matter.
\item \DIFdelbegin \DIFdel{Structure-resolved $n_f$ estimates yield settling velocities 2--10$\times$ lower }\DIFdelend \DIFaddbegin \DIFadd{Revised fractal dimensions indicate lower settling rates }\DIFaddend than conventional aggregated estimates\DIFdelbegin \DIFdel{, implying longer transport distances}\DIFdelend .
\end{keypoints}


\begin{abstract}
Mud builds coastal landscapes and \DIFdelbegin \DIFdel{hosts pollutants and particulate }\DIFdelend \DIFaddbegin \DIFadd{carries pollutants and }\DIFaddend carbon, yet predicting its transport is difficult because mud aggregates into flocs \DIFdelbegin \DIFdel{, which control settling rates in riverine}\DIFdelend \DIFaddbegin \DIFadd{that control its settling rate in river}\DIFaddend , estuarine, and marine \DIFdelbegin \DIFdel{environments}\DIFdelend \DIFaddbegin \DIFadd{settings}\DIFaddend . Flocs are typically characterized using fractal theory with a constant fractal dimension ($n_f = 2.0$\DIFdelbegin \DIFdel{--$2.2$), inferred from aggregated }\DIFdelend \DIFaddbegin \DIFadd{), which is inferred from aggregating }\DIFaddend floc size--settling velocity data\DIFdelbegin \DIFdel{despite their }\DIFdelend \DIFaddbegin \DIFadd{, despite a diversity of }\DIFaddend complex fractal structures. We conducted experiments on freshwater flocs under varied shear stress and organic \DIFdelbegin \DIFdel{matter concentrations }\DIFdelend \DIFaddbegin \DIFadd{concentration }\DIFaddend to characterize this diversity. \DIFdelbegin \DIFdel{The fractal dimension $n_f$ was calculated from coupled measurements of floc size and settling rate using particle image tracking. }\DIFdelend Rather than the conventional approach of aggregating \DIFdelbegin \DIFdel{all }\DIFdelend settling data to infer \DIFdelbegin \DIFdel{a single value of }\DIFdelend $n_f$, we grouped measurements by the two-dimensional \DIFdelbegin \DIFdel{fractal }\DIFdelend \DIFaddbegin \DIFadd{(2D) box-counting }\DIFaddend dimension derived from \DIFdelbegin \DIFdel{box counting}\DIFdelend \DIFaddbegin \DIFadd{particle tracking}\DIFaddend , thereby partially resolving structural variability. This stratified \DIFdelbegin \DIFdel{approach }\DIFdelend \DIFaddbegin \DIFadd{method }\DIFaddend reveals that $n_f$ increases with shear velocity \DIFdelbegin \DIFdel{(from 1.77 to 2.12) }\DIFdelend and decreases with organic matter content\DIFdelbegin \DIFdel{(from 1.84 to 1.63)}\DIFdelend , trends that are obscured when data are aggregated. \DIFdelbegin \DIFdel{Across our conditions, aggregation biases inferred $n_f$ by $\sim$5--10\%, which propagates nonlinearly into settling velocities and predicted transport distances. The relationship between 2D and 3D fractal dimensions is best described by a power-law conversion ($n_f = 0.81\,(n_f^{(2\mathrm{D})})^{1.81}$; }%DIFDELCMD < \citeA{wang2022fractal}%%%
\DIFdel{), consistent with the sub-Euclidean space filling of porous aggregates. Although conventional bulk fits to our data yield }\DIFdelend \DIFaddbegin \DIFadd{The stratified analysis resolves }\DIFaddend $n_f$ \DIFdelbegin \DIFdel{values near the canonical range, we show that these estimates are biased due to correlations among settling velocity, floc size, and fractal dimension, analogous to Simpson's Paradox}\DIFdelend \DIFaddbegin \DIFadd{variations of 0.2--0.35 units across experimental conditions, whereas the conventional aggregated approach reduces this range by roughly a factor of five. The inferred primary particle diameter $d_p$ co-varies with these conditions, and the concurrent variation of $n_f$, $d_p$, and floc size introduces competing effects on settling velocity}\DIFaddend . Our approach \DIFdelbegin \DIFdel{reduces this bias, }\DIFdelend yields lower $n_f$ \DIFdelbegin \DIFdel{, and implies }\DIFdelend \DIFaddbegin \DIFadd{values (1.6--2.1) than the canonical range, which propagates nonlinearly into settling rates, implying }\DIFaddend that mud flocs are more porous \DIFdelbegin \DIFdel{, }\DIFdelend \DIFaddbegin \DIFadd{and }\DIFaddend settle more slowly \DIFdelbegin \DIFdel{, and travel farther than previous models predict. This finding suggests that freshwater }\DIFdelend \DIFaddbegin \DIFadd{than expected. We also find that the 2D-to-3D $n_f$ relationship can be described by a nonlinear model, consistent with empirical conversions for synthetic aggregates. Overall, our results indicate that }\DIFaddend mud retention in floodplains and deltas is less efficient than current floc models \DIFdelbegin \DIFdel{assume---at least under }\DIFdelend \DIFaddbegin \DIFadd{assume, at least in }\DIFaddend the freshwater conditions \DIFdelbegin \DIFdel{analyzed here---with }\DIFdelend \DIFaddbegin \DIFadd{we analyzed, with }\DIFaddend implications for sediment budgets, contaminant dispersal, and coastal landscape evolution.
\DIFdelbegin %DIFDELCMD < 

%DIFDELCMD < %%%
\DIFdelend \end{abstract}

\section*{Plain Language Summary}
Mud particles in rivers tend to clump together into clusters called flocs. How quickly these flocs sink determines where mud is deposited, shaping the growth of river deltas, floodplains, and coastlines. Predicting this settling is challenging because flocs vary widely in structure: some are dense and compact, while others are porous and loosely bound. Most studies assume that all flocs have a single, uniform structure, which oversimplifies their behavior. In our experiments, we generated flocs in freshwater under controlled conditions with different flow strengths and amounts of organic matter. Using high-resolution imaging, we measured both their structure and settling speed. We found that faster flows tend to produce \DIFaddbegin \DIFadd{smaller but }\DIFaddend more compact flocs\DIFdelbegin \DIFdel{that settle more quickly}\DIFdelend , while organic matter encourages the formation of more open structures that settle more slowly. By grouping flocs according to their measured structural complexity, we revealed trends that are hidden when data are \DIFdelbegin \DIFdel{grouped }\DIFdelend \DIFaddbegin \DIFadd{averaged }\DIFaddend together. This new approach reduces bias in estimating settling rates and improves predictions of how mud moves through rivers and into estuaries and oceans, where salinity may further modify \DIFdelbegin \DIFdel{physical clumping }\DIFdelend \DIFaddbegin \DIFadd{aggregation }\DIFaddend beyond the freshwater trends isolated here.

\section{Introduction}

Fine-grained sediment, or mud, defined by particles smaller than 62.5~\textmu m, dominates fluvial sediment loads and represents the \DIFdelbegin \DIFdel{primary }\DIFdelend \DIFaddbegin \DIFadd{dominant }\DIFaddend material deposited in lowland and deltaic environments \cite{healy2002muddy, kranenburg1994fractal, lamb2020mud, nghiem2026flocculated}. The dynamics of mud transport and deposition therefore exert \DIFaddbegin \DIFadd{a }\DIFaddend first-order control on the long-term evolution of deltas, floodplains, and estuaries \cite{nghiem2022mechanistic, winterwerp2002flocculation}, as well as on water quality, contaminant fluxes, and the success of restoration interventions \cite{droppo1997freshwater, esposito2017efficient, kemp2016enhancing, liss1996floc}. Mud also plays a central role in the global carbon cycle by providing reactive mineral surfaces for organic carbon adsorption and long-term preservation \cite{bianchi2024anthropogenic, blair2012fate}.

The geomorphic and geochemical importance of mud necessitates a mechanistic framework for particle transport and settling to enable predictive modeling. Settling velocity is a key parameter in sediment-transport equations and is commonly estimated using Stokes\DIFdelbegin \DIFdel{' }\DIFdelend \DIFaddbegin \DIFadd{’ }\DIFaddend law, which gives the terminal settling velocity ($w_s$) of smooth, impermeable spheres at low Reynolds numbers \cite{bird2006transport, stokes1851effect}:
\begin{equation}
w_s = \DIFdelbegin \DIFdel{\frac{g\,(\rho_p - \rho_w)\,d^2}{18\,\rho_w\,\nu}
    }\DIFdelend \DIFaddbegin \DIFadd{\frac{g\,(\rho_p - \rho_f)\,d^2}{18\,\rho_f\,\nu}
}\DIFaddend \label{eq:stokes}
\end{equation}
where $g = 9.81\ \mathrm{m\,s^{-2}}$ is gravitational acceleration, $\rho_p\ (\mathrm{kg\,m^{-3}})$ and \DIFdelbegin \DIFdel{$\rho_w\ (\mathrm{kg\,m^{-3}})$ }\DIFdelend \DIFaddbegin \DIFadd{$\rho_f\ (\mathrm{kg\,m^{-3}})$ }\DIFaddend are the densities of the \DIFdelbegin \DIFdel{primary particle and water, respectively}\DIFdelend \DIFaddbegin \DIFadd{particle and fluid}\DIFaddend , $d\ (\mathrm{m})$ is particle diameter, and $\nu\ (\mathrm{m^2\,s^{-1}})$ is the kinematic viscosity of the \DIFdelbegin \DIFdel{water}\DIFdelend \DIFaddbegin \DIFadd{fluid}\DIFaddend .

However, when applied to fine-grained sediment, Stokes\DIFdelbegin \DIFdel{' }\DIFdelend \DIFaddbegin \DIFadd{’ }\DIFaddend law yields predictions at odds with field observations. Consider a single \DIFdelbegin \DIFdel{1~}%DIFDELCMD < \textmu %%%
\DIFdel{m primary mineral grain (e.g., quartz or clay, $\rho_p \approx 2650\ \mathrm{kg\,m^{-3}}$) }\DIFdelend \DIFaddbegin \DIFadd{$1~\text{\textmu m}$ mud particle }\DIFaddend in a $5$\DIFdelbegin \DIFdel{\,}\DIFdelend \DIFaddbegin \DIFadd{~}\DIFaddend m-deep river flowing at \DIFdelbegin \DIFdel{$1\,\mathrm{m\,s}^{-1}$}\DIFdelend \DIFaddbegin \DIFadd{$1~\mathrm{m\,s^{-1}}$}\DIFaddend . Using Stokes\DIFdelbegin \DIFdel{' }\DIFdelend \DIFaddbegin \DIFadd{’ }\DIFaddend law produces a settling velocity of \DIFdelbegin \DIFdel{approximately $9\times 10^{-7}\,\mathrm{m\,s}^{-1}$}\DIFdelend \DIFaddbegin \DIFadd{just $9\times10^{-8}~\mathrm{m\,s^{-1}}$}\DIFaddend , implying a deposition time \DIFdelbegin \DIFdel{of about 64 }\DIFdelend \DIFaddbegin \DIFadd{greater than $600$~}\DIFaddend days and a transport distance exceeding \DIFdelbegin \DIFdel{5,500\,}\DIFdelend \DIFaddbegin \DIFadd{$55{,}000$~}\DIFaddend km downstream. This hypothetical scenario conflicts sharply with field observations of mud deposition \DIFaddbegin \DIFadd{in fluvial and coastal settings }\DIFaddend \cite{lamb2020mud, mason2019differential, nicholas1996significance}\DIFdelbegin \DIFdel{. Field studies document that lowland rivers commonly build levees through overbank flows, with many levees composed primarily of mud and clay contents of approximately $50\%$ }%DIFDELCMD < \cite{aalto2008spatial}%%%
\DIFdel{. Furthermore, deposition rates along levees typically decrease exponentially with distance from the channel, indicating transport lengths on the order of tens to hundreds of meters rather than thousands of kilometers }%DIFDELCMD < \cite{hajek2012simplified}%%%
\DIFdelend .

This discrepancy indicates that mud settles far more rapidly than predicted for isolated particles, \DIFdelbegin \DIFdel{implying the presence of }\DIFdelend \DIFaddbegin \DIFadd{pointing to }\DIFaddend additional physical processes. The deviation arises because suspended mud particles flocculate, forming larger aggregates \DIFdelbegin \DIFdel{that are }\DIFdelend composed of individual primary particles \DIFdelbegin \DIFdel{(representing the constituent mineral grains) and }\DIFdelend that settle much faster than their \DIFdelbegin \DIFdel{isolated }\DIFdelend \DIFaddbegin \DIFadd{individual }\DIFaddend constituents \cite{droppo1997freshwater, johnson1996settling, kranenburg1994fractal}. These flocculated aggregates, or flocs, are highly porous rather than dense, impermeable spheres, so their effective density decreases with increasing size \cite{kranenburg1994fractal, Meakin1992Aggregation}. Because flocs exhibit approximately self-similar fractal organization \cite{family2012kinetics, jullien1987aggregation}, their hierarchical structure can be described by a \DIFdelbegin \DIFdel{standard }\DIFdelend fractal mass--size scaling \cite{Meakin1992Aggregation}:
\begin{equation}
N \sim \left( \frac{d_f}{d_p} \right)^{n_f},
\label{eq:fractaltheory}
\end{equation}
where $N$ is the number of primary particles (dimensionless), $d_f$ is the floc diameter (m), $d_p$ is the primary particle diameter (m), and $n_f$ is the \DIFdelbegin \DIFdel{3D }\DIFdelend fractal dimension (dimensionless).

Building on Eq.~\eqref{eq:fractaltheory}, the fraction of the floc volume \DIFdelbegin \DIFdel{actually }\DIFdelend occupied by solid grains is the ratio of solid volume to the enclosing floc volume. Because the solid volume scales with $N$ times the monomer volume while the enclosure scales with $d_f^{3}$, the solid volume fraction \DIFdelbegin \DIFdel{$\phi_s$ }\DIFdelend \DIFaddbegin \DIFadd{$\phi$ }\DIFaddend decreases with size approximately as \DIFdelbegin \DIFdel{$(d_f/d_p)^{\,n_f-3}$}\DIFdelend \DIFaddbegin \DIFadd{$(d_f/d_p)^{n_f-3}$}\DIFaddend . For a porous aggregate whose pores are filled by ambient water, the floc\DIFdelbegin \DIFdel{'}\DIFdelend \DIFaddbegin \DIFadd{’}\DIFaddend s submerged specific gravity equals this solid fraction times the \DIFdelbegin \DIFdel{primary particle'}\DIFdelend \DIFaddbegin \DIFadd{solid’}\DIFaddend s submerged specific gravity. Denoting submerged specific gravity by \DIFdelbegin \DIFdel{$R=(\rho_p-\rho_w)/\rho_w$}\DIFdelend \DIFaddbegin \DIFadd{$R = (\rho - \rho_w)/\rho_w$}\DIFaddend , this yields:
\begin{equation}
R_f = R\DIFdelbegin \DIFdel{_p }\DIFdelend \DIFaddbegin \DIFadd{_s }\DIFaddend \left( \frac{d_f}{d_p} \right)\DIFdelbegin \DIFdel{^{n_f - 3}}\DIFdelend \DIFaddbegin \DIFadd{^{n_f-3}}\DIFaddend ,
\label{eq:eq3}
\end{equation}
where $R_f$ and \DIFdelbegin \DIFdel{$R_p$ }\DIFdelend \DIFaddbegin \DIFadd{$R_s$ }\DIFaddend are the submerged specific \DIFdelbegin \DIFdel{gravities }\DIFdelend \DIFaddbegin \DIFadd{gravity }\DIFaddend of the floc and primary particles, respectively (both dimensionless), \DIFdelbegin \DIFdel{$\rho_p$ is primary }\DIFdelend \DIFaddbegin \DIFadd{$\rho$ is }\DIFaddend particle density (kg\,m$^{-3}$), and $\rho_w$ is water density (kg\,m$^{-3}$). Since a dense solid sphere has \DIFdelbegin \DIFdel{a }\DIFdelend volume that scales with the cube of its diameter, its fractal dimension is \DIFdelbegin \DIFdel{$n_f=3$}\DIFdelend \DIFaddbegin \DIFadd{$n_f = 3$}\DIFaddend , and Eq.~\eqref{eq:eq3} reduces to \DIFdelbegin \DIFdel{$R_f=R_p$}\DIFdelend \DIFaddbegin \DIFadd{$R_f = R_s$}\DIFaddend . Naturally formed flocs are more porous with $1 \le n_f < 3$, yielding \DIFdelbegin \DIFdel{$R_f<R_p$}\DIFdelend \DIFaddbegin \DIFadd{$R_f < R_s$}\DIFaddend . Accordingly, floc effective (excess) density decreases with floc size. For a fractal aggregate, the solids mass scales as $M_f \propto d_f^{n_f}$ whereas the envelope volume scales as $V_f \propto d_f^{3}$, implying \DIFdelbegin \DIFdel{the excess density, $\Delta \rho_f$, scales as }\DIFdelend $\Delta \rho_f \propto M_f / V_f \propto d_f^{n_f-3}$ \cite{kranenburg1994fractal}.

Substituting Eq.~\eqref{eq:eq3} into Eq.~\eqref{eq:stokes} gives an explicit expression for the settling velocity of a fractal floc \cite{strom2011explicit, winterwerp1998simple}:
\begin{equation}
w_s = \DIFdelbegin \DIFdel{\frac{g\,R_p}{b_1\,\nu\,d_p^{\,n_f - 3}}\; }\DIFdelend \DIFaddbegin \DIFadd{\frac{g\,R_s}{b_1\,\nu\,d_p^{\,n_f-3}}\, }\DIFaddend d_f\DIFdelbegin \DIFdel{^{\,n_f - 1}}\DIFdelend \DIFaddbegin \DIFadd{^{\,n_f-1}}\DIFaddend ,
\label{eq:floc_settling}
\end{equation}
where $w_s$ is the settling velocity (m\,s$^{-1}$)\DIFaddbegin \DIFadd{, $g$ is gravitational acceleration (m\,s$^{-2}$), $\nu$ is kinematic viscosity (m$^{2}$\,s$^{-1}$), }\DIFaddend and $b_1$ is a dimensionless shape factor (with \DIFdelbegin \DIFdel{$b_1=18$ }\DIFdelend \DIFaddbegin \DIFadd{$b_1 = 18$ }\DIFaddend for smooth spheres).

In practice, applying Eq.~\eqref{eq:floc_settling} requires an appropriate fractal dimension $n_f$, which captures how floc structure modifies settling velocity. Experimental and field studies often assume \DIFdelbegin \DIFdel{$n_f \approx 2.0$--$2.2$ }\DIFdelend \DIFaddbegin \DIFadd{$n_f \approx 2$ }\DIFaddend from empirical $w_s$\DIFdelbegin \DIFdel{--}\DIFdelend \DIFaddbegin \DIFadd{–}\DIFaddend $d_f$ relationships, values reported primarily for saline estuarine/coastal flocs where electrochemical interactions promote more compact aggregates; freshwater flocs more often show lower $n_f$ ($\lesssim 2$), with overlap depending on turbulence and composition \DIFdelbegin %DIFDELCMD < \cite{winterwerp1998simple,winterwerp2006heuristic}%%%
\DIFdelend \DIFaddbegin \cite{winterwerp1998simple, winterwerp2006heuristic}\DIFaddend . With all other parameters held fixed, Eq.~\eqref{eq:floc_settling} gives \DIFdelbegin \DIFdel{$w_s \propto d_f^{n_f-1}$ }\DIFdelend \DIFaddbegin \DIFadd{$w_s \propto d_f^{\,n_f-1}$ }\DIFaddend \cite{johnson1996settling, li1989fractal}. Measurements of $w_s$ and $d_f$ across a floc population therefore allow $n_f$ to be inferred from the \DIFdelbegin \DIFdel{log--log }\DIFdelend \DIFaddbegin \DIFadd{log–log }\DIFaddend slope of the $w_s$\DIFdelbegin \DIFdel{--}\DIFdelend \DIFaddbegin \DIFadd{–}\DIFaddend $d_f$ relation \cite{kranenburg1994fractal, strom2011explicit, winterwerp1998simple, winterwerp2002flocculation, winterwerp2004introduction, winterwerp2006heuristic}. In general, this inference assumes that floc permeability is independent of floc size. If permeability instead varies systematically with $d_f$, the observed scaling relation may reflect both fractal structure and permeability effects \cite{strom2011explicit, nghiem2026flocculated}.

Using a constant population-wide value for the fractal dimension is convenient but restrictive. \citeA{khelifa2006models} showed that treating $n_f$ as constant fails to reproduce the observed spread in settling velocities and argued that $n_f$ can decrease with floc size, motivating models that prescribe a function $n_f(d_f)$, often as a power law, to represent within-population variability. In practice, however, most experimental and image-based field studies measure only floc size $d_f$ and settling velocity $w_s$ \cite{smith2015image, nghiem2024testing}, so the fractal dimension is not independently observed and functions for $n_f$ cannot be directly tested. Moreover, models for $n_f(d_f)$ implicitly assume a unique fractal dimension for a given floc size \cite{khelifa2006models}, whereas natural floc populations can contain multiple flocs with comparable $d_f$ but distinct internal structure, and therefore distinct $n_f$ \cite{maggi2007variable}. As a result, standard \DIFdelbegin \DIFdel{log--log }\DIFdelend \DIFaddbegin \DIFadd{log–log }\DIFaddend regressions of $w_s$ on $d_f$ \DIFdelbegin \DIFdel{implicitly }\DIFdelend treat structurally heterogeneous populations as homogeneous, and \DIFdelbegin \DIFdel{data aggregation of }\DIFdelend \DIFaddbegin \DIFadd{pooling }\DIFaddend such subpopulations can misattribute between-group differences to size scaling itself, producing biased estimates of $n_f$.

To illustrate this effect, consider a population comprising three subgroups of flocs, each defined by a distinct fractal dimension ($n_f = 1.5$, $2.0$, and $2.5$) but formed from primary particles of the same diameter ($d_p = 10^{-5}$ \DIFdelbegin \DIFdel{~}\DIFdelend m). According to Eq.~\eqref{eq:floc_settling}, the slope of the \DIFdelbegin \DIFdel{log--log }\DIFdelend \DIFaddbegin \DIFadd{log–log }\DIFaddend relationship between settling velocity and floc diameter for each subgroup is $n_f - 1$. When analyzed separately, these subgroups yield distinct \DIFdelbegin \DIFdel{log--log }\DIFdelend \DIFaddbegin \DIFadd{log–log }\DIFaddend slopes corresponding to $n_f = 1.5$, $2.0$, and $2.5$, respectively (Fig.~\ref{fig:1}). However, if these floc subgroups are combined into a single group and floc diameter ($d_f$) is correlated with the 3D fractal dimension ($n_f$), then fitting a single power-law relationship to all points can yield a substantially skewed and misleading exponent. In Fig.~\ref{fig:1}, the resulting \DIFdelbegin \DIFdel{log--log }\DIFdelend \DIFaddbegin \DIFadd{log–log }\DIFaddend slope is $2$, implying \DIFdelbegin \DIFdel{an effective }\DIFdelend $n_f = 3$, which corresponds to the settling behavior of smooth, impermeable spheres, not flocs. This example parallels Simpson\DIFdelbegin \DIFdel{'}\DIFdelend \DIFaddbegin \DIFadd{’}\DIFaddend s Paradox \cite{simpson1951interpretation}, a statistical phenomenon in which relationships observed within individual subgroups are reversed, masked, or distorted \DIFdelbegin \DIFdel{upon data aggregation}\DIFdelend \DIFaddbegin \DIFadd{when data are pooled}\DIFaddend . In the context of floc settling, \DIFdelbegin \DIFdel{data aggregation of }\DIFdelend \DIFaddbegin \DIFadd{pooling }\DIFaddend structurally distinct subpopulations can produce a global $w_s$\DIFdelbegin \DIFdel{--}\DIFdelend \DIFaddbegin \DIFadd{–}\DIFaddend $d_f$ slope that implies an apparent $n_f$ inconsistent with the $n_f$ of any subgroup, so the inferred ``representative'' fractal dimension is a regression artifact rather than a physical descriptor.
\DIFdelbegin %DIFDELCMD < 

%DIFDELCMD < %%%
\DIFdelend \begin{figure}[htbp]
    \centering
    \DIFdelbeginFL \DIFdelFL{\includegraphics[width=0.8\textwidth]{figures/1.pdf}
    }\DIFdelendFL \DIFaddbeginFL \DIFaddFL{\includegraphics[width=0.8\textwidth]{figures/fig1.pdf}
    }\DIFaddendFL \caption{
    Synthetic data illustrating Simpson's Paradox in floc settling. Three floc subgroups, each with a specific 3D fractal dimension ($n_f = 1.5,\ 2.0,\ \text{and}\ 2.5$), are shown. By Eq.~\eqref{eq:floc_settling}, the slope of the log--log relationship between settling velocity ($w_s$) and floc diameter ($d_f$) is $n_f - 1$. Fitting each subgroup individually recovers slopes of 0.5, 1.0, and 1.5, as expected. In contrast, \DIFdelbeginFL \DIFdelFL{data aggregation of }\DIFdelendFL \DIFaddbeginFL \DIFaddFL{pooling }\DIFaddendFL all points produces a single power-law fit (red line) with a misleading slope of 2, equivalent to $n_f = 3$, characteristic of solid spheres and not flocs. All trend lines intersect at the primary particle diameter ($d_f = d_p = 10^{-5}\,\mathrm{m}$), where settling velocity is independent of fractal dimension.
    }
    \label{fig:1}
\end{figure}

To estimate fractal dimension independently of settling velocity and floc diameter, studies have employed image-based techniques that analyze in situ 2D projections of flocs \cite{jarvis2005measuring}. These approaches typically calculate a two-dimensional (2D) fractal dimension from high-resolution images, then infer the corresponding three-dimensional (3D) structure via empirical or simulation-based relationships. A common method relates the perimeter--area scaling of projected flocs to a 2D fractal dimension, using computationally generated aggregates to derive a mapping between 2D and 3D fractal dimension \cite{lee2004prediction, maggi2004method, tang2015reconstructing}. However, the relationship between 2D and 3D fractal dimension is derived from synthetic aggregates, which may not replicate the aggregation dynamics or structural variability of natural flocs. Projection-based estimates therefore risk systematic bias and, more importantly, do not recover the 3D fractal dimension that directly controls settling behavior.

Although flocculation has been widely characterized, how within-population structural heterogeneity affects settling remains uncertain, particularly under natural particle-size distributions and aggregation pathways. The central difficulty is that \DIFdelbegin \DIFdel{$w_s - d_f$ }\DIFdelend \DIFaddbegin \DIFadd{$w_s$–$d_f$ }\DIFaddend regressions alone do not condition on structure, so \DIFdelbegin \DIFdel{data aggregation of }\DIFdelend \DIFaddbegin \DIFadd{pooling }\DIFaddend distinct subpopulations can produce an apparent relation and inferred $n_f$ that is biased by \DIFaddbegin \DIFadd{data-pooling (}\DIFaddend Simpson-type\DIFaddbegin \DIFadd{) }\DIFaddend effects. To address this bias, we use images from experiments to stratify the floc population by the \DIFdelbegin \DIFdel{two-dimensional (}\DIFdelend 2D \DIFdelbegin \DIFdel{) fractal dimension from }\DIFdelend box-counting \DIFdelbegin \DIFdel{, }\DIFdelend \DIFaddbegin \DIFadd{dimension }\DIFaddend and within each subgroup, regress settling velocity $w_s$ against floc diameter $d_f$ to estimate the hydrodynamically relevant \DIFdelbegin \DIFdel{three-dimensional (}\DIFdelend 3D \DIFdelbegin \DIFdel{) }\DIFdelend fractal dimension $n_f$. This conditioning step provides a principled way to separate structural heterogeneity from size scaling without prescribing a population-level closure such as $n_f(d_f)$: that is, $n_f$ is inferred hydrodynamically within approximately homogeneous structural classes. The assumption is that flocs with similar 2D \DIFdelbegin \DIFdel{fractal }\DIFdelend \DIFaddbegin \DIFadd{box-counting }\DIFaddend dimension have similar 3D fractal dimension, which we validate herein.

We evaluate this framework in controlled laboratory experiments designed to simulate flocs in rivers with systematic variation in shear velocity \DIFdelbegin \DIFdel{($u_*$) }\DIFdelend and organic matter concentration. High-resolution in situ imaging enabled quantification of floc morphology and settling trajectories. Specifically, we pursued three objectives: (1) demonstrate the existence of structurally distinct subgroups within the floc population, each characterized by a unique fractal dimension; (2) \DIFdelbegin \DIFdel{evaluate the relationship between }\DIFdelend \DIFaddbegin \DIFadd{relate }\DIFaddend 2D \DIFdelbegin \DIFdel{and 3D fractal dimension}\DIFdelend \DIFaddbegin \DIFadd{image–derived fractal dimensions to settling-inferred 3D dimensions}\DIFaddend , and assess this relationship against theoretical predictions; and (3) estimate the effective primary particle size ($d_p$) and quantify how the hydrodynamic shape factor ($b_1$) varies as a function of $n_f$. All experiments were conducted in freshwater (salinity $\approx 0$), so the results isolate the roles of shear and organic binding and are not intended as a direct constraint for saline estuaries where salt-driven cohesion can further modify flocculation.

\section{Methodology}

\subsection{Experimental Setup}

We performed flocculation experiments in a custom mixing tank designed to provide controlled conditions (Fig.~\ref{fig:2}a). The tank comprises a cylindrical PVC pipe, 20~cm in diameter and 60~cm tall, filled with water to a depth of 40~cm. \DIFaddbegin \DIFadd{The tank was filled with fresh tap water, used without filtration or chemical modification, and experiments were conducted at room temperature. }\DIFaddend A motor-driven rotating lid, mounted above the open top and submerged just below the free surface, sets the desired shear. \DIFaddbegin \DIFadd{For each experiment, the tank was filled with water and the motor was started to establish the target shear. Particles and organic matter (described below) were then added gradually, and the tank was mixed for 30~min to allow flocs to form under sustained shear. The motor was subsequently stopped to initiate settling, and video recording began immediately and continued for the first 30~min.
}

\DIFaddend Shear velocity, $u_*$, was empirically calibrated for each run by converting motor revolutions per minute to $u_*$ using a relationship derived from previous sediment-entrainment trials in \citeA{douglas2025mud}. We define $u_* = (\tau/\rho_w)^{1/2}$, where $\tau$ is a characteristic \DIFdelbegin \DIFdel{turbulent fluid shear stressand $\rho_w$ is the density of water}\DIFdelend \DIFaddbegin \DIFadd{bed shear stress}\DIFaddend , and we treat $u_*$ as \DIFdelbegin \DIFdel{an effective, run-averaged }\DIFdelend \DIFaddbegin \DIFadd{a characteristic }\DIFaddend measure of turbulent shear intensity because the \DIFdelbegin \DIFdel{rotating-lid-driven flow field is }\DIFdelend \DIFaddbegin \DIFadd{rotating lid-driven flow field and boundary stresses are }\DIFaddend spatially heterogeneous.

During the mixing phase, the flow was fully turbulent. This is quantified by the friction Reynolds number $Re_{\tau} = u_* h / \nu$, where $h$ is the water depth and $\nu$ is the kinematic viscosity of freshwater. For $h = 0.40$~m and $\nu \approx 1.0 \times 10^{-6}\ \mathrm{m^2\,s^{-1}}$ at room temperature, the imposed shear velocities $u_* = 0.02$\DIFdelbegin \DIFdel{--}\DIFdelend \DIFaddbegin \DIFadd{–}\DIFaddend $0.04$~m~s$^{-1}$ correspond to $Re_{\tau} \approx (8$\DIFdelbegin \DIFdel{--}\DIFdelend \DIFaddbegin \DIFadd{–}\DIFaddend $16)\times10^{3}$, well within the fully turbulent regime for wall-bounded shear flows \DIFdelbegin %DIFDELCMD < \cite{Pope2000,Schlichting2000}%%%
\DIFdelend \DIFaddbegin \cite{Pope2000, Schlichting2000}\DIFaddend . Across the flocculation experiments, we imposed $u_* = 0.02$\DIFdelbegin \DIFdel{--}\DIFdelend \DIFaddbegin \DIFadd{–}\DIFaddend $0.04$~m~s$^{-1}$ (equivalent to $\tau \approx 0.4$\DIFdelbegin \DIFdel{--1.6}\DIFdelend \DIFaddbegin \DIFadd{–$1.6$}\DIFaddend ~Pa), spanning low-to-moderate shear conditions typical of suspension-dominated river channels, sufficient for sand entrainment but below those required to mobilize gravels \cite{julien2018river}.

\DIFaddbegin \DIFadd{All image-based measurements reported here were collected during the settling phase. }\DIFaddend Suspended flocs were imaged through a 3~\DIFdelbegin \DIFdel{mm-wide }\DIFdelend \DIFaddbegin \DIFadd{mm-thick }\DIFaddend flow-through slit built into the tank wall at the observation level. A DSLR fitted with a $5\times$ microscope objective was aligned with the slit to record flocs as they crossed the narrow optical path, enhancing image clarity and increasing the number of focused particles captured by the camera, following methods similar to \citeA{osborn2021flocarazi}.

\begin{figure}[htbp]
    \centering
    \DIFdelbeginFL \DIFdelFL{\includegraphics[width=1\textwidth]{figures/2.pdf}
    }\DIFdelendFL \DIFaddbeginFL \DIFaddFL{\includegraphics[width=1\textwidth]{figures/fig2.pdf}
    }\DIFaddendFL \caption{
    (a) Schematic of the custom flocculation apparatus. A cylindrical PVC tank (20~cm diameter, 60~cm height) is filled to 40~cm depth and stirred with a motor-driven \DIFaddbeginFL \DIFaddFL{rotating }\DIFaddendFL lid just below the surface. \DIFdelbeginFL \DIFdelFL{Shear velocity is calibrated for each run. }\DIFdelendFL An observation slit (3~mm wide) enables \DIFdelbeginFL \textit{\DIFdelFL{in situ}} %DIFAUXCMD
\DIFdelendFL \DIFaddbeginFL \DIFaddFL{in situ }\DIFaddendFL imaging of suspended flocs with a DSLR camera and \DIFdelbeginFL \DIFdelFL{5$\times$ }\DIFdelendFL \DIFaddbeginFL \DIFaddFL{$5\times$ }\DIFaddendFL microscope objective aligned for high-resolution focus. \DIFdelbeginFL \DIFdelFL{Key components and dimensions are annotated.
    }\DIFdelendFL (b) Grain-size distributions of the dispersed input sediment measured by \DIFaddbeginFL \DIFaddFL{the }\DIFaddendFL Mastersizer from \citeA{nghiem2025flocculation}. Curves show pure silica sand (SiO$_2$), montmorillonite clay, and their 50:50 volume-weighted mixture.}
    \label{fig:2}
\end{figure}

\subsection{Experimental Runs}

We conducted eight experimental runs in total. The first two were calibration runs using laboratory spheres and silica grains, which do not flocculate, to verify the accuracy of our settling-velocity measurements and provide a baseline for comparing non-flocculating particles to flocs under identical experimental conditions. The remaining six runs systematically varied shear velocity \DIFdelbegin \DIFdel{($u_*$) }\DIFdelend and organic matter (OM) content, given that previous studies have shown \DIFaddbegin \DIFadd{that }\DIFaddend turbulent shear limits floc size by promoting breakup, while organic matter enhances \DIFdelbegin \DIFdel{physical }\DIFdelend aggregation by binding particles \cite{nghiem2022mechanistic}. By independently controlling these variables, we aimed to isolate their respective roles in floc formation and settling.

\DIFdelbegin \DIFdel{For each experiment, the tank was filled with water and the motor was started to establish the target shear. Particles and organic matter were then added gradually, and the tank was mixed for 30~min to allow flocs to form under sustained shear. The motor was subsequently stopped to initiate settling, and video recording began immediately and continued for the first 30~min. All measurements of floc size, settling velocity, and fractal dimension reported in this study were collected during this settling phase, not during mixing. }\DIFdelend Each floc experiment used 10~g of solids: 5~g silica (SiO$_2$) and 5~g montmorillonite clay. Organic content was varied in three runs by adding guar gum, a \DIFdelbegin \DIFdel{plant-derived polysaccharide used as a proxy for natural organic }\DIFdelend \DIFaddbegin \DIFadd{polysaccharide proxy for plant-derived extracellular }\DIFaddend polymers, at 1, 2, or 3~wt\% relative to the total mass of solids, following the organic-loading framework of \citeA{zeichner2021early}. This range was selected to (i) remain within the lower end of organic fractions reported for natural river suspended sediment (\DIFdelbegin \DIFdel{on the order of 1--6}\DIFdelend \DIFaddbegin \DIFadd{order 1–6}\DIFaddend \% and up to \DIFdelbegin \DIFdel{\(\sim\)}\DIFdelend \DIFaddbegin \DIFadd{$\sim$}\DIFaddend 20\% in compiled field datasets), and (ii) avoid polymer-dominated regimes while still producing a resolvable change in flocculation dynamics. In \citeA{zeichner2021early}, polymer-to-clay ratios were chosen explicitly to span modern river organic-matter concentrations, and their experiments show that even \DIFdelbegin \DIFdel{1--2}\DIFdelend \DIFaddbegin \DIFadd{1–2}\DIFaddend ~wt\% guar gum produces a strong flocculation response. Guar gum is a controlled proxy for natural extracellular polymeric substances and does not capture the full chemical complexity of field organic matter. Full experimental conditions, including organic-matter fraction and imposed $u_*$, are listed in Table~\ref{tab:exp_conditions}.

\begin{table}[htbp]
\caption{Summary of experimental conditions.\label{tab:exp_conditions}}
\centering
\begin{tabular}{lccc}
\hline
Exp Num & Experiment Type & OM (wt\%) & Shear velocity (m/s) \\
\hline
1 & Lab Sphere     & NA    & NA    \\
2 & Silica         & NA    & NA    \\
3 & Floc Shear 1   & 2     & 0.02  \\
4 & Floc Shear 2   & 2     & 0.03  \\
5 & Floc Shear 3   & 2     & 0.04  \\
6 & Floc OM 1     & 1     & 0.04  \\
7 & Floc OM 2     & 2     & 0.04  \\
8 & Floc OM 3     & 3     & 0.04  \\
\hline
\multicolumn{4}{l}{NA: Not applicable.}
\end{tabular}
\end{table}

Grain-size distributions for the silica and montmorillonite were measured using a Malvern Mastersizer 3000 (Fig.~\ref{fig:2}b), following the procedure of \citeA{nghiem2025flocculation}. In brief, samples were dispersed in deionized water, sonicated to break up aggregates, and analyzed by laser diffraction to obtain particle-size distributions. Together, these two minerals span nearly three orders of magnitude in diameter, \DIFdelbegin \DIFdel{yielding a volume-weighted size distribution }\DIFdelend \DIFaddbegin \DIFadd{with primary particles ranging }\DIFaddend from $2 \times 10^{-7}$ to $2 \times 10^{-4}$~m \DIFaddbegin \DIFadd{(percent finer by volume)}\DIFaddend . For the mixed-sediment runs, a composite distribution was calculated as the 50:50 volume-weighted average of the two minerals.

\subsection{Image Processing\DIFdelbegin \DIFdel{Methods for Floc Settling Velocity and Diameter}\DIFdelend }

The floc tracking analysis began with processing every frame of the recorded video sequence, acquired at a constant frame interval \DIFdelbegin \DIFdel{$\Delta t = 1/30~\mathrm{s}$, as verified }\DIFdelend \DIFaddbegin \DIFadd{$\Delta t = 1/30$~s, as read }\DIFaddend from the video metadata\DIFaddbegin \DIFadd{, so that the sampling frequency equals the acquisition rate}\DIFaddend . Each frame was converted to an eight-bit grayscale image using OpenCV, and a time-averaged background image (computed by averaging a block of 1000 consecutive frames) was subtracted from each image to isolate transient features such as moving flocs (Fig.~\ref{fig:3}a).

\begin{figure}[htbp]
    \centering
    \DIFdelbeginFL \DIFdelFL{\includegraphics[width=1\textwidth]{figures/3.pdf}
    }\DIFdelendFL \DIFaddbeginFL \DIFaddFL{\includegraphics[width=1\textwidth]{figures/fig3.pdf}
    }\DIFaddendFL \caption{
    Image processing workflow used for floc analysis. (a) Example of a background-subtracted image, highlighting transient objects such as flocs. (b) Detected and contoured particles passing focus and sharpness thresholds (outlined in red) with raw velocity vectors overlaid. (c) Local background flow field calculated by PIV from tracer particles. (d) Example of a tracked floc with both its raw and corrected (background-subtracted) settling velocity vectors indicated.
    }
    \label{fig:3}
\end{figure}

Particle boundaries in each image were detected from intensity gradients using a global threshold of 200 applied to the normalized residual image. To retain only sharp and well-focused flocs, candidate objects were required to exceed a Laplacian variance threshold of 5 (measuring focus/sharpness) and to have an intensity standard deviation greater than 5 within their contour \cite{pertuz2013analysis}. These specific threshold values were optimized for our experimental conditions and image resolution \DIFdelbegin \DIFdel{, }\DIFdelend and may vary between different setups. Only particles meeting both criteria were accepted for further analysis and are highlighted in red in Fig.~\ref{fig:3}b. This objective\DIFaddbegin \DIFadd{, }\DIFaddend per-frame filtering mitigates focus-related sampling bias by excluding blurred detections rather than allowing size-dependent blur to enter the floc statistics. For each accepted object, we recorded the planform area in image pixels $A_{\mathrm{px}}$ (later converted to an equivalent diameter after spatial calibration) and the centroid coordinates $(x_c,y_c)$ \DIFaddbegin \DIFadd{(}\DIFaddend for subsequent velocity estimation\DIFaddbegin \DIFadd{)}\DIFaddend .

To express diameters and velocities in physical units, we calibrated pixel size using a precision stage micrometer imaged under the same optical conditions as the experiments. Comparing the pixel spacing between micrometer markings with their true separations yielded a scale factor of \(0.45~\mu\mathrm{m\,px^{-1}}\). We applied this factor to all subsequent analyses to convert particle dimensions and inter-frame displacements from pixels to micrometers; velocities were then obtained from these displacements using the recorded frame interval.

Settling velocities were obtained by frame-to-frame tracking using a conservative nearest-neighbor association. For each detection in frame \(t\), we searched within a 50-pixel radius in frame \(t{+}1\) for candidate matches. We scored each candidate \(j\) using two equally weighted, normalized differences: the equivalent concave-hull diameter \(d_{\mathrm{ch}}\) (px) and the Laplacian-variance sharpness \(L\) (arb. units). The score was defined as \(S_j=\tfrac{1}{2}\big(|d_{\mathrm{ch},j}-d_{\mathrm{ch},t}|/d_{\mathrm{ch},t}+|L_j-L_t|/L_t\big)\). We selected the candidate with the smallest \(S_j\), provided that both relative differences were \(\le 0.10\); otherwise, the track was terminated to avoid spurious links from rapid shape or focus changes. The vertical velocity for that interval was initially computed as \(\Delta y\,s/\Delta t\), where \(\Delta y\) is the vertical centroid shift (px, positive downward), \(s\) is the spatial calibration (\(0.45~\mu\mathrm{m\,px^{-1}}\)), and \(\Delta t\) is the frame interval (s). Raw velocities computed in \(\mu\mathrm{m\,s^{-1}}\) were subsequently converted to \(\mathrm{m\,s^{-1}}\) to match the physical units required by Eq.~\eqref{eq:stokes} and Eq.~\eqref{eq:floc_settling}.

Frame-to-frame particle displacements yield raw \DIFdelbegin \DIFdel{two-dimensional }\DIFdelend \DIFaddbegin \DIFadd{2D }\DIFaddend velocity vectors that combine gravitational settling with background flow in the imaging plane. To isolate the settling signal, we first computed a local 2D flow field using OpenPIV \cite{liberzon2021openpiv} for each frame pair. Following \citeA{smith2015image}, the images were digitally filtered to retain only tracers smaller than $10~\mu\mathrm{m}$ that behave as passive flow followers (Fig.~\ref{fig:3}c). For each tracked floc, we then subtracted the corresponding local flow vector---bilinearly interpolated at the floc centroid---from the floc's raw trajectory vector. This vector differencing produces a background-corrected motion (Fig.~\ref{fig:3}d, gray arrow for raw, white arrow for corrected), from which we report the vertical component as the settling velocity \(w_s\). All settling velocities in this study refer to this corrected, background-subtracted quantity.

After velocity correction, floc diameters were derived from each particle's projected area. Two common image-based metrics are used to estimate diameter: the Feret diameter (maximum caliper length across the projection) and the equivalent circular diameter (ECD), defined as the diameter of a circle with the same projected area \cite{jarvis2005measuring, smith2015image}. The Feret diameter is sensitive to orientation and outliers, often overestimating size for irregular or elongated flocs. By contrast, the area-equivalent diameter---especially when computed from a concave hull---can better represent irregular shapes by de-emphasizing extreme protrusions, though it may underrepresent thin filaments. Accordingly, we adopt the concave-hull ECD (with $\alpha = 0.05$, the concave-hull radius parameter that controls the allowed boundary indentation), corresponding to the projected-area--equivalent diameter of \citeA{smith2015image}, while acknowledging that alternative definitions are reasonable \cite{jarvis2005measuring}. Particles were segmented with binary masks, boundaries were traced using a concave-hull algorithm, the enclosed area was measured, and the diameter of a circle of equal area was computed and used as \(d_f\) in all subsequent log--log analyses. Because $d_f$ is defined geometrically from the 2D projection, this step does not assume constant floc density; any density changes associated with size or organic content enter only through the measured settling velocities.

\subsection{Calibration and Settling \DIFdelbegin \DIFdel{Behavior }\DIFdelend Analysis}

We calibrated our setup with Experiments~1 and~2, which measured the settling of laboratory spheres (50~\textmu m) and natural silica grains (Fig.~\ref{fig:4}a). Settling velocities were fitted with Stokes' law (Eq.~\eqref{eq:stokes}), treating the shape factor $b_1$ as a free parameter. For the spheres, the best-fit value $b_1 = 17.7$ is consistent with the theoretical value of $18$ for smooth spheres (Fig.~\ref{fig:4}b), which lies within the confidence interval of the fit. This agreement indicates that the experimental setup accurately reproduces Stokes' settling behavior for smooth spheres. In contrast, silica grains settled more slowly, yielding $b_1 = 29.2$, a drag increase attributable to their angular, non-spherical shapes; \citeA{ferguson2004simple} reported a comparable $b_1 \approx 24$ for natural sediments of similar morphology.

Compared with silica, flocs settled more slowly and showed substantially greater dispersion in velocity at a given diameter (Fig.~\ref{fig:4}b). This reduced settling reflects the lower effective density of flocs arising from their highly porous, irregular internal structure \cite{kranenburg1994fractal, winterwerp1998simple}. Consequently, the floc data are not described by the standard Stokes formulation (Eq.~\eqref{eq:stokes}) and instead require the floc-specific relation (Eq.~\eqref{eq:floc_settling}), which introduces the 3D fractal dimension $n_f$ and the primary particle diameter $d_p$ as additional parameters.

\begin{figure}[htbp]
    \centering
    \DIFdelbeginFL \DIFdelFL{\includegraphics[width=1\textwidth]{figures/4.pdf}
    }\DIFdelendFL \DIFaddbeginFL \DIFaddFL{\includegraphics[width=1\textwidth]{figures/fig4.pdf}
    }\DIFaddendFL \caption{
    (a) Images of a laboratory sphere (Exp.~1), silica grain (Exp.~2), and a representative floc (Exp.~5). The 30~$\mu$m scale bar applies to all. (b) Log--log plot of settling velocity ($w_s$) vs.\ diameter ($d_f$) for spheres (green circles, Exp.~1), silica (red circles, Exp.~2), and aggregated flocs (blue circles, Exps.~3--8). Sphere and silica fits use Stokes' law (Eq.~\eqref{eq:stokes}), yielding shape factors ($b_1$) of 17.7 and 29.2. The floc data fit the fractal settling model ($w_s \propto d_f^{n_f - 1}$) with a slope of 0.90, giving an apparent fractal dimension $n_f = 1.90$. (c) Sensitivity of the fractal settling model (Eq.~\eqref{eq:floc_settling}) to assumed primary particle diameter ($d_p$). The floc data are compared to theoretical predictions using $n_f = 1.9$ for a primary particle size ($d_p$) of $10^{-6}$~m (dashed line) and $10^{-5}$~m (solid red line). This highlights how the choice of $d_p$ controls the required model fit, a key uncertainty in the conventional approach.
    }
    \label{fig:4}
\end{figure}

However, directly calculating $n_f$ by fitting Eq.~\eqref{eq:floc_settling} to the floc data is problematic because the primary-particle diameter, $d_p$, is unknown. Although our experiments have known distributions of input particle sizes (Fig.~\ref{fig:2}b), we do not know which particle sizes are actually incorporated into any given floc. Even if we could measure the primary particle sizes directly, the fractal theory necessitates specifying only a single value for a characteristic primary particle size, and it is unclear how to define this value from a distribution. Moreover, because $n_f$ appears both as an exponent on $d_f$ and within the pre-factor of Eq.~\eqref{eq:floc_settling} (via $d_p^{-(n_f-3)}$), a conventional regression that fits the model to the data for a prescribed $d_p$ (a one-parameter fit) forces the inferred $n_f$ to mathematically compensate for the assumed pre-factor. Consequently, the assumed value for $d_p$ exerts a strong influence on the inferred fractal dimension $n_f$ (Fig.~\ref{fig:4}c): assuming $d_p = 1\times 10^{-6}$~m yields $n_f = 2.77$; $d_p = 1\times 10^{-5}$~m gives $n_f = 2.10$; and $d_p = 1\times 10^{-4}$~m produces $n_f = 3.00$.

To avoid explicitly solving for the unknown primary-particle size $d_p$, we use the logarithmic form of Eq.~\eqref{eq:floc_settling}: $\log(w_s) = (n_f - 1)\log(d_f) + C$. Here, the intercept $C$ lumps together $d_p$, the shape factor $b_1$, and other constants. Under the assumption that \(d_p\) and \(b_1\) are approximately constant across the fitted population, the slope in \(\log(w_s)\)--\(\log(d_f)\) space yields \(n_f = \text{slope} + 1\) without requiring prior knowledge of \(d_p\). For the full floc population, the fitted slope is \(0.90\), giving \(n_f = 1.90\) (Fig.~\ref{fig:4}b). However, if structural heterogeneity causes \(n_f\) to covary with size, or if effective \(d_p\) varies across the population, a bulk regression can exhibit Simpson-type data-aggregation bias. We therefore stratify flocs by their \DIFdelbegin \textit{\DIFdel{in situ}}%DIFAUXCMD
\DIFdelend \DIFaddbegin \DIFadd{in situ}\DIFaddend , image-based fractal dimension and fit separate regressions within each stratum; the resulting slopes provide subgroup-specific estimates of \(n_f\) and reveal systematic dependence on structure.

\subsection{\DIFdelbegin \DIFdel{Estimating Fractal Dimension of Flocs from }\DIFdelend \DIFaddbegin \DIFadd{Image-based }\DIFaddend 2D \DIFdelbegin \DIFdel{Projections}\DIFdelend \DIFaddbegin \DIFadd{box-counting dimension}\DIFaddend }

Fractal theory relates the fractal dimension to the primary particle diameter, the number of primary particles, and the floc diameter (Eq.~\eqref{eq:fractaltheory}). In natural systems, however, neither $d_p$ nor $N$ is constant or directly measurable during experiments. As a result, numerous methods have been developed to estimate floc fractal dimensions from image-based properties. These approaches are inherently \DIFdelbegin \DIFdel{two-dimensional}\DIFdelend \DIFaddbegin \DIFadd{2D}\DIFaddend , yielding an apparent fractal dimension \DIFdelbegin \DIFdel{$n_f^{(2\mathrm{D})}$ }\DIFdelend that may differ from the hydrodynamically relevant \DIFdelbegin \DIFdel{three-dimensional }\DIFdelend \DIFaddbegin \DIFadd{3D }\DIFaddend value. Throughout this paper, we denote the 3D fractal dimension \DIFdelbegin \DIFdel{simply }\DIFdelend \DIFaddbegin \DIFadd{inferred from settling velocity scaling }\DIFaddend as $n_f$, and use \DIFdelbegin \DIFdel{$n_f^{(2\mathrm{D})}$ when referring specifically to }\DIFdelend \DIFaddbegin \DIFadd{$n_f^{(\mathrm{2D})}$ for }\DIFaddend the 2D \DIFdelbegin \DIFdel{image-based measure}\DIFdelend \DIFaddbegin \DIFadd{box-counting dimension measured from images. Where the specific 2D method matters (Methodology), we distinguish the box-counting dimension $n_{f,\,\mathrm{box}}^{(\mathrm{2D})}$ from the perimeter--area dimension $n_{f,\,\mathrm{perimeter}}^{(\mathrm{2D})}$; in the Results, all 2D values refer to box-counting and are written simply as $n_f^{(\mathrm{2D})}$}\DIFaddend .

The perimeter--area method has previously been used to estimate \DIFdelbegin \DIFdel{the }\DIFdelend \DIFaddbegin \DIFadd{a }\DIFaddend 2D fractal dimension of flocs~\cite{lee2004prediction, maggi2004method, tang2015reconstructing}. This method relies on the empirical scaling relation:
\begin{equation}
P \propto A\DIFdelbegin \DIFdel{^{n_{f,\,\text{perim}}^{\text{(2D)}} / 2}
}\DIFdelend \DIFaddbegin \DIFadd{^{n_{f,\,\mathrm{perimeter}}^{(\mathrm{2D})} / 2}
}\DIFaddend \end{equation}
where $P$ is the measured perimeter and $A$ is the projected area of the particle in the 2D image. In practice, \DIFdelbegin \DIFdel{$n_{f,\,\text{perim}}^{\text{(2D)}}$ }\DIFdelend \DIFaddbegin \DIFadd{$n_{f,\,\mathrm{perimeter}}^{(\mathrm{2D})}$ }\DIFaddend is estimated from the log--log slope of $P$ versus $A$. This expression is derived by analogy to the Euclidean relation $P = k A^{1/2}$, which holds for smooth (non-fractal) boundaries with \DIFdelbegin \DIFdel{$n_f^{(2\mathrm{D})} = 1$}\DIFdelend \DIFaddbegin \DIFadd{$n_{f,\,\mathrm{perimeter}}^{(\mathrm{2D})} = 1$}\DIFaddend . The underlying assumption is that, for a fractal boundary, the exponent in this scaling law approximates the true boundary fractal dimension \DIFdelbegin \DIFdel{$n_f^{(2\mathrm{D})}$}\DIFdelend \DIFaddbegin \DIFadd{$n_{f,\,\mathrm{perimeter}}^{(\mathrm{2D})}$}\DIFaddend . However, both mathematical analysis and geometric constructions (e.g., the Koch curve) demonstrate that this assumption breaks down for truly fractal sets~\cite{cheng1995perimeter,frame_mandelbrot_neger_2025}. For such boundaries, the measured perimeter is strongly dependent on the image resolution\DIFdelbegin \DIFdel{(the `yardstick' length)}\DIFdelend , making a simple $P$--$A$ scaling ill-defined without explicitly accounting for the measurement scale. As a result, perimeter--area scaling does not reliably recover the true fractal dimension from 2D projections.

Box-counting, in contrast, is an appropriate method for fractal-dimension estimation that does not rely on strict self-similarity and is computationally straightforward \cite{foroutan1999advances}. A key advantage of box-counting is that it does not require knowledge of the primary particle size ($d_p$) or the number of constituent particles ($N$), making it particularly well-suited for natural flocs where these values are unknown or variable~\cite{bellouti1997flocs, spencer2022quantification}. The box-counting dimension \DIFdelbegin \DIFdel{$n_{f,\text{box}}^{\text{(2D)}}$ }\DIFdelend \DIFaddbegin \DIFadd{$n_{f,\,\mathrm{box}}^{(\mathrm{2D})}$ }\DIFaddend is determined by measuring how the number of occupied boxes $N_{\text{box}}(\epsilon)$ at a given scale $\epsilon$ varies with scale~\cite{mandelbrot1967long}:
\begin{equation}
N_{\text{box}}(\epsilon) \propto \epsilon\DIFdelbegin \DIFdel{^{-n_{f,\text{box}}^{\text{(2D)}}}
}\DIFdelend \DIFaddbegin \DIFadd{^{-n_{f,\,\mathrm{box}}^{(\mathrm{2D})}}
}\DIFaddend \end{equation}

For each floc mask, we computed $N_{\text{box}}(\epsilon)$ using square boxes with side lengths $\epsilon$ spanning powers of two in pixel units (\(\epsilon=1,2,4,\dots\)), and we estimated the dimension from the slope of a least-squares fit over the portion of the log--log curve that remains approximately linear and spans at least one order of magnitude in scale. When analyzing images of objects with finite extent from \DIFdelbegin \textit{\DIFdel{in situ}} %DIFAUXCMD
\DIFdelend \DIFaddbegin \DIFadd{in situ }\DIFaddend imaging, finite-size effects and image pixelation make the estimated fractal dimension sensitive to the available scaling range. For a digital image of $W \times W$ pixels, the practical range of box sizes ($\epsilon$) spans from 1 to $W$ pixels. Reliable estimation of \DIFdelbegin \DIFdel{$n_{f,\text{box}}^{\text{(2D)}}$ }\DIFdelend \DIFaddbegin \DIFadd{$n_{f,\,\mathrm{box}}^{(\mathrm{2D})}$ }\DIFaddend requires a sufficiently broad range of box sizes to establish a stable linear trend on a log--log plot of $N_{\text{box}}(\epsilon)$ versus $1/\epsilon$. If the image is too small, the limited scaling range can yield unstable values of \DIFdelbegin \DIFdel{$n_{f,\text{box}}^{\text{(2D)}}$}\DIFdelend \DIFaddbegin \DIFadd{$n_{f,\,\mathrm{box}}^{(\mathrm{2D})}$}\DIFaddend . For example, a $20 \times 20$ pixel image provides only about 1.3 orders of magnitude in scale, which is insufficient for robust fitting and produces scale-dependent estimates (Fig.~\ref{fig:5}a). In contrast, a $100 \times 100$ pixel image spans two full orders of magnitude and typically provides 6--7 usable data points for regression; beyond this image size, the estimated \DIFdelbegin \DIFdel{$n_{f,\text{box}}^{\text{(2D)}}$ }\DIFdelend \DIFaddbegin \DIFadd{$n_{f,\,\mathrm{box}}^{(\mathrm{2D})}$ }\DIFaddend saturates, indicating a practical threshold for resolution-independent estimation (Fig.~\ref{fig:5}a). Accordingly, all analyzed flocs were extracted using $100 \times 100$ pixel windows centered on each particle. Fig.~\ref{fig:5}b shows examples for two experimental flocs at high image resolution, illustrating the application of the box-counting method to real data. In this study, all reported values of the 2D \DIFdelbegin \DIFdel{fractal dimension ($n_{f}^{\text{(2D)}}$}\DIFdelend \DIFaddbegin \DIFadd{box-counting dimension ($n_{f,\,\mathrm{box}}^{(\mathrm{2D})}$}\DIFaddend ) in the results and discussion sections refer exclusively to measurements obtained by the box-counting method.

\begin{figure}[htbp]
    \centering
    \DIFdelbeginFL \DIFdelFL{\includegraphics[width=1\textwidth]{figures/5.pdf}
    }\DIFdelendFL \DIFaddbeginFL \DIFaddFL{\includegraphics[width=1\textwidth]{figures/fig5.pdf}
    }\DIFaddendFL \caption{
    (a) Illustration of resolution dependence in box-counting dimension estimation. Top: A perfect circle embedded in $20 \times 20$ and $200 \times 200$ pixel grids shows that coarse grids systematically underestimate the fractal dimension. Bottom: Plot of measured 2D \DIFdelbeginFL \DIFdelFL{fractal }\DIFdelendFL \DIFaddbeginFL \DIFaddFL{box-counting }\DIFaddendFL dimension \DIFdelbeginFL \DIFdelFL{$n_f^{(2\mathrm{D})}$ }\DIFdelendFL \DIFaddbeginFL \DIFaddFL{$n_{f,\,\mathrm{box}}^{(\mathrm{2D})}$ }\DIFaddendFL versus image width $W$, highlighting a resolution-dependent regime at low $W$ and a stable region at high $W$. The transition scale ($W \approx 100$ pixels) was identified by repeating the box-counting calculation across increasing image sizes and selecting the smallest $W$ beyond which further increases produced negligible changes in the estimated dimension.
    (b) Example box-counting dimension estimates for two experimental flocs imaged at $>100 \times 100$ pixel resolution.
    }
    \label{fig:5}
\end{figure}

\subsection{\DIFdelbegin \DIFdel{Estimating }\DIFdelend Primary Particle Diameter ($d_p$)\DIFdelbegin \DIFdel{from Settling Data}\DIFdelend }

In most \DIFdelbegin \textit{\DIFdel{in~situ}} %DIFAUXCMD
\DIFdelend \DIFaddbegin \DIFadd{in situ }\DIFaddend experiments, the primary particle diameter \(d_p\) is assumed from the input sediment-size distribution and hypothesized floc-formation conditions because \(d_p\) is rarely resolvable directly. This practice can be uncertain when the effective primary scale differs from the feed distribution or varies across aggregation pathways. Here we identify \(d_p\) empirically by leveraging multiple floc subpopulations with distinct fractal dimensions \(n_f\). In the experimental application, we treat $d_p$ as an experiment-level parameter common to all $n_f$-stratified subpopulations. As implied by Eq.~\eqref{eq:floc_settling}, \(d_p\) does not affect the slope of the \(\log(w_s)\)--\(\log(d_f)\) relation (the slope is \(n_f-1\)), but it fixes a common intersection of the subgroup regressions. This is because at \(d_f=d_p\), the object is a primary particle (not a floc), so substituting \(d_f=d_p\) into Eq.~\eqref{eq:floc_settling} yields \DIFdelbegin \DIFdel{\(w_s=(g\,R_p/b_1\nu)\,d_p^{2}\)}\DIFdelend \DIFaddbegin \DIFadd{\(w_s=(g\,R_s/b_1\nu)\,d_p^{2}\)}\DIFaddend , which is independent of \(n_f\). Consequently, regressions from any number of \(n_f\)-stratified subpopulations intersect at \DIFdelbegin \DIFdel{\((d_f, w_s)=(d_p,\,(g\,R_p/b_1\nu)\,d_p^{2})\)}\DIFdelend \DIFaddbegin \DIFadd{\((d_f, w_s)=(d_p,\,(g\,R_s/b_1\nu)\,d_p^{2})\)}\DIFaddend . We refer to this common intersection diameter as $d_\text{intersect}$. In our synthetic example (Fig.~\ref{fig:1}), three subgroups with \(n_f \approx 1.5,\,2.0,\) and \(2.5\) produce regression lines that converge at \(d_\text{intersect} \approx 10^{-5}\,\mathrm{m}\), because the synthetic data were generated using this prescribed primary-particle diameter across varying fractal dimensions, thereby illustrating the method.


\DIFdelbegin \DIFdel{For typical observed flocs, the measured sizes imply that each }\DIFdelend \DIFaddbegin \DIFadd{A }\DIFaddend floc contains many primary particles. For example, with \DIFdelbegin \DIFdel{$d_f/d_p \sim 10$--$100$ }\DIFdelend \DIFaddbegin \DIFadd{$d_f \sim 10^{-4}$~m, $d_p \sim 10^{-5}$~m, }\DIFaddend and $n_f \approx 2$, Eq.~\eqref{eq:fractaltheory} gives \DIFdelbegin \DIFdel{$N \sim 10^{2}$--}\DIFdelend \DIFaddbegin \DIFadd{$N \sim (d_f/d_p)^{n_f} \sim 10^{2}$–}\DIFaddend $10^{4}$. This means the settling measurements reflect the collective behavior of aggregates built from a common sediment pool, and the inferred $d_p$ should be interpreted as an effective primary scale for the flocs resolved by the imaging and tracking. In the synthetic example, the subgroup regressions intersect exactly at a single point because the data are generated under the idealized assumption that all primary particles share a single diameter $d_p$. In natural sediments, primary particles span a distribution of sizes, so different flocs can be built from \DIFdelbegin \DIFdel{different effective primary scales }\DIFdelend \DIFaddbegin \DIFadd{differently sized particles }\DIFaddend and the subgroup regressions need not converge to a unique intersection. We therefore treat the intersection as distributed rather than singular and use a Monte Carlo procedure to generate an ensemble of intersection points, which we interpret as a distribution of effective $d_p$ values. We then compare this inferred $d_p$ distribution to the input sediment-size distribution to assess consistency and to quantify how the sediment\DIFdelbegin \DIFdel{'}\DIFdelend \DIFaddbegin \DIFadd{’}\DIFaddend s primary-size variability propagates into the settling-based estimate of $d_p$.

The intersection method assumes a common shape factor $b_1$ across fractal subgroups. Because drag depends on aggregate structure, $b_1$ may vary with $n_f$, biasing the inferred intersection diameter. Allowing for subgroup-dependent drag, the relationship between the observed intersection diameter $d_\text{intersect}$ and the true primary particle diameter $d_p$ for two subpopulations $j$ and $k$ is
\begin{equation}
d_\text{intersect} = d_p \left( \frac{b_{1,j}}{b_{1,k}} \right)^{1/(n_{f,j} - n_{f,k})},
\label{eq:dpintersection}
\end{equation}
which reduces to
\begin{equation}
d_\text{intersect} = d_p
\label{eq:dpintersectionideal}
\end{equation}
when $b_{1,j}=b_{1,k}$. In this limit, all subgroup regressions intersect exactly at the primary particle diameter.

\section{Results}

\subsection{\DIFdelbegin \DIFdel{Bulk }\DIFdelend Floc \DIFdelbegin \DIFdel{Size and Settling Velocity}\DIFdelend \DIFaddbegin \DIFadd{Bulk Properties}\DIFaddend }
\DIFdelbegin %DIFDELCMD < 

%DIFDELCMD < %%%
\DIFdel{Before stratifying by fractal structure, we }\DIFdelend \DIFaddbegin \label{sec:bulk}
\DIFadd{We }\DIFaddend first report the bulk distributions of floc diameter ($d_f$) and settling velocity ($w_s$) across all experimental conditions (\DIFdelbegin \DIFdel{Table~\ref{tab:bulk_trends}). These summary statistics characterize the first-order response of the floc population to the imposed shear and organic matter forcing.
}%DIFDELCMD < 

%DIFDELCMD < \begin{table}[htbp]
%DIFDELCMD < %%%
%DIFDELCMD < \caption{%
{%DIFAUXCMD
\DIFdelFL{Median floc diameter and settling velocity with interquartile ranges (IQR) for each experimental condition.}%DIFDELCMD < \label{tab:bulk_trends}%%%
}
%DIFAUXCMD
%DIFDELCMD < \centering
%DIFDELCMD < \begin{tabular}{lcccccc}
%DIFDELCMD < \hline
%DIFDELCMD < %%%
\DIFdelFL{Experiment }%DIFDELCMD < & %%%
\DIFdelFL{$u_*$ (m/s) }%DIFDELCMD < & %%%
\DIFdelFL{OM (wt\%)}%DIFDELCMD < & %%%
\DIFdelFL{$N$ }%DIFDELCMD < & %%%
\DIFdelendFL \DIFaddbeginFL \DIFaddFL{Fig.~\ref{fig:6}). Median }\DIFaddendFL $d_f$ \DIFdelbeginFL \DIFdelFL{(}%DIFDELCMD < \textmu %%%
\DIFdelFL{m) }%DIFDELCMD < [%%%
\DIFdelFL{IQR}%DIFDELCMD < ] & %%%
\DIFdelFL{$w_s$ (mm/s)}%DIFDELCMD < [%%%
\DIFdelFL{IQR}%DIFDELCMD < ] \\
%DIFDELCMD < \hline
%DIFDELCMD < %%%
\DIFdelFL{LOWSHEAR  }%DIFDELCMD < & %%%
\DIFdelFL{0.02 }%DIFDELCMD < & %%%
\DIFdelFL{2 }%DIFDELCMD < & %%%
\DIFdelFL{2429 }%DIFDELCMD < & %%%
\DIFdelFL{43.2 }%DIFDELCMD < [%%%
\DIFdelFL{38.1--57.9}%DIFDELCMD < ] & %%%
\DIFdelFL{0.30 }%DIFDELCMD < [%%%
\DIFdelFL{0.21--0.44}%DIFDELCMD < ] \\
%DIFDELCMD < %%%
\DIFdelFL{MIDSHEAR  }%DIFDELCMD < & %%%
\DIFdelFL{0.03 }%DIFDELCMD < & %%%
\DIFdelFL{2 }%DIFDELCMD < & %%%
\DIFdelFL{2201 }%DIFDELCMD < & %%%
\DIFdelFL{28.0 }%DIFDELCMD < [%%%
\DIFdelFL{24.6--35.9}%DIFDELCMD < ] & %%%
\DIFdelFL{0.39 }%DIFDELCMD < [%%%
\DIFdelFL{0.28--0.55}%DIFDELCMD < ] \\
%DIFDELCMD < %%%
\DIFdelFL{HIGHSHEAR }%DIFDELCMD < & %%%
\DIFdelFL{0.04 }%DIFDELCMD < & %%%
\DIFdelFL{2 }%DIFDELCMD < & %%%
\DIFdelFL{4603 }%DIFDELCMD < & %%%
\DIFdelFL{26.8 }%DIFDELCMD < [%%%
\DIFdelFL{23.1--32.4}%DIFDELCMD < ] & %%%
\DIFdelFL{0.32 }%DIFDELCMD < [%%%
\DIFdelFL{0.24--0.46}%DIFDELCMD < ] \\
%DIFDELCMD < %%%
\DIFdelFL{OM1       }%DIFDELCMD < & %%%
\DIFdelFL{0.04 }%DIFDELCMD < & %%%
\DIFdelendFL \DIFaddbeginFL \DIFaddFL{decreases monotonically with increasing shear velocity (Fig.~\ref{fig:6}a), consistent with stronger turbulence breaking apart aggregates. In contrast, $d_f$ responds weakly and non-monotonically to organic matter content, increasing slightly from }\DIFaddendFL 1\DIFdelbeginFL %DIFDELCMD < & %%%
\DIFdelFL{1058 }%DIFDELCMD < & %%%
\DIFdelFL{42.0 }%DIFDELCMD < [%%%
\DIFdelFL{35.7--50.7}%DIFDELCMD < ] & %%%
\DIFdelFL{0.78 }%DIFDELCMD < [%%%
\DIFdelFL{0.63--1.00}%DIFDELCMD < ] \\
%DIFDELCMD < %%%
\DIFdelFL{OM2       }%DIFDELCMD < & %%%
\DIFdelFL{0.04 }%DIFDELCMD < & %%%
\DIFdelendFL \DIFaddbeginFL \DIFaddFL{\% to }\DIFaddendFL 2\DIFdelbeginFL %DIFDELCMD < & %%%
\DIFdelFL{943  }%DIFDELCMD < & %%%
\DIFdelFL{44.6 }%DIFDELCMD < [%%%
\DIFdelFL{36.8--55.6}%DIFDELCMD < ] & %%%
\DIFdelFL{0.67 }%DIFDELCMD < [%%%
\DIFdelFL{0.52--0.88}%DIFDELCMD < ] \\
%DIFDELCMD < %%%
\DIFdelFL{OM3       }%DIFDELCMD < & %%%
\DIFdelFL{0.04 }%DIFDELCMD < & %%%
\DIFdelendFL \DIFaddbeginFL \DIFaddFL{\% OM before declining at }\DIFaddendFL 3\DIFdelbeginFL %DIFDELCMD < & %%%
\DIFdelFL{1453 }%DIFDELCMD < & %%%
\DIFdelFL{38.8 }%DIFDELCMD < [%%%
\DIFdelFL{32.0--50.5}%DIFDELCMD < ] & %%%
\DIFdelFL{0.52 }%DIFDELCMD < [%%%
\DIFdelFL{0.39--0.73}%DIFDELCMD < ] \\
%DIFDELCMD < \hline
%DIFDELCMD < \end{tabular}
%DIFDELCMD < \end{table}
%DIFDELCMD < 

%DIFDELCMD < %%%
\DIFdel{In the shear-velocity series (Experiments 3--5; constant OM = 2~wt\%) , median floc diameter decreases from 43.2~}%DIFDELCMD < \textmu %%%
\DIFdel{m at $u_* = 0.02$}\DIFdelend \DIFaddbegin \DIFadd{\% OM (Fig.~\ref{fig:6}b). Median settling velocity $w_s$ exhibits a non-monotonic dependence on shear velocity, peaking at $u_* = 0.03$}\DIFaddend ~m\DIFdelbegin \DIFdel{\,s$^{-1}$ to 26.8~}%DIFDELCMD < \textmu %%%
\DIFdel{m at $u_* = 0.04$~m \,s$^{-1}$, consistent with shear-driven breakup limiting maximum floc size }%DIFDELCMD < \cite{winterwerp2002flocculation}%%%
\DIFdel{. Despite the decrease in size, median settling velocity does not decrease proportionally: it is0.30~mm\,s$^{-1}$ at low shear, rises to 0.39~mm\,s}\DIFdelend \DIFaddbegin \DIFadd{\,s}\DIFaddend $^{-1}$ \DIFdelbegin \DIFdel{at mid shear, and returns to 0.32~mm\,s$^{-1}$ at high shear. This non-monotonic pattern reflects the competing effects of decreasing floc size and increasing floc compactness (higher $n_f$)with shear; smaller but denser flocs can maintain comparable settling velocities}\DIFdelend \DIFaddbegin \DIFadd{rather than decreasing steadily (Fig.~\ref{fig:6}c). With increasing organic matter, $w_s$ decreases monotonically (Fig.~\ref{fig:6}d).
}

\begin{figure}[htbp]
    \centering
    \DIFaddFL{\includegraphics[width=0.8\textwidth]{figures/fig6.pdf}
    }\caption{\DIFaddFL{Median floc diameter ($d_f$) and settling velocity ($w_s$) with interquartile ranges for varying shear velocity (a, c) and organic matter content (b, d). Dashed lines show normalized model predictions: in panels (a, b), the }\citeA{nghiem2022mechanistic} \DIFaddFL{equilibrium floc diameter trend $d_{f,\text{eq}} = f(\eta,\theta)$, which is independent of $n_f$; in panels (c, d), the base model with $n_f = 2$ (red dashed; Eq.~\eqref{eq:floc_settling} with $n_f=2$) and the modified model with measured $n_f$ and $d_p$ (blue dashed; Eq.~\eqref{eq:floc_settling} with condition-specific $n_f$ and $d_p$; see Section~\ref{sec:reconciling}). Each model trend is normalized by a single least-squares constant.}}
    \label{fig:6}
\end{figure}


\DIFadd{To interpret these trends mechanistically, we compare our observations against the semi-empirical freshwater flocculation framework of }\citeA{nghiem2022mechanistic}\DIFadd{. That model predicts both the equilibrium floc diameter $d_f$ and the corresponding settling velocity $w_s$ as functions of two environmental inputs: the Kolmogorov microscale $\eta$ (which sets the smallest turbulent eddy that can break flocs) and the fractional organic surface coverage $\theta$ (which controls binding strength). Evaluating this model against our data requires mapping our experimental control variables---imposed shear velocity $u_*$ and organic matter weight fraction---to $\eta$ and $\theta$}\DIFaddend .

In \DIFdelbegin \DIFdel{the organic-matter series (Experiments 6--8; constant $u_* = 0.04$~}\DIFdelend \DIFaddbegin \DIFadd{applying their model to river field data, }\citeA{nghiem2022mechanistic} \DIFadd{used an open-channel relation where turbulent dissipation is generated solely by bed friction. Our experimental apparatus is a top-driven cylindrical tank in which stationary sidewall friction acts as an additional dominant energy sink. To calculate the volume-averaged turbulent dissipation rate $\bar{\varepsilon}_{\text{tank}}$, we evaluate the steady-state power balance }\cite{camp1943velocity}\DIFadd{, equating the power injected by the rotating lid to the power dissipated by friction on all tank boundaries (bed and sidewalls). Assuming the wall shear stress scales with the calibrated bed shear velocity ($\tau_w \approx \rho_w u_*^2$) }\cite{Pope2000, Schlichting2000} \DIFadd{and substituting the bulk slip velocity $U_{\text{bulk}} \approx u_* \sqrt{2/C_f}$, the mean dissipation rate is:
}\begin{equation}
    \DIFadd{\bar{\varepsilon}_{\text{tank}} = \frac{u_*^3}{\sqrt{C_f/2}} \left( \frac{1}{h} + \frac{2}{R} \right),
    \label{eq:eps_tank}
}\end{equation}
\DIFadd{where $h = 0.40$~}\DIFaddend m \DIFaddbegin \DIFadd{is the water depth, $R = 0.10$~m is the tank radius, and $C_f \approx 0.005$ is the skin friction coefficient for fully turbulent swirling flow }\cite{daily1960chamber, Schlichting2000}\DIFadd{. The Kolmogorov microscale is then $\eta = (\nu^3 / \bar{\varepsilon}_{\text{tank}})^{1/4}$. Because of the extensive sidewall area, $\bar{\varepsilon}_{\text{tank}}$ is substantially larger than open-channel predictions for an equivalent $u_*$, so the absolute value of $\eta$ is smaller. However, $\bar{\varepsilon}_{\text{tank}} \propto u_*^3$ regardless of tank geometry, so the functional scaling is identical to an open channel: $\eta \propto u_*^{-3/4}$.
}

\citeA{nghiem2022mechanistic} \DIFadd{parameterized organic binding using field measurements of particulate organic carbon (OC), requiring an intermediate stoichiometric assumption (cellulose composition) to convert OC to total organic matter (OM). Because our experiments use guar gum---a pure polysaccharide proxy for natural extracellular polymeric substances---we bypass this conversion and use the experimental weight percentage of OM directly. Following Eq.~19 of }\citeA{nghiem2022mechanistic}\DIFadd{, the fractional surface coverage is:
}\begin{equation}
    \DIFadd{\theta = \frac{(\text{\%OM}/100)\,(\rho_s / \rho_{\text{OM}})\, d_p^3}{(d_p + \delta)^3 - d_p^3},
    \label{eq:theta}
}\end{equation}
\DIFadd{where $\rho_s = 2650$~kg}\DIFaddend \,\DIFdelbegin \DIFdel{s$^{-1}$), median settling velocity decreases systematically from 0.78~mm\,s$^{-1}$ at 1~wt\% OM to 0.52~mm}\DIFdelend \DIFaddbegin \DIFadd{m$^{-3}$ is the sediment grain density, $\rho_{\text{OM}} = 1000$~kg}\DIFaddend \,\DIFdelbegin \DIFdel{s$^{-1}$ at 3~wt\% OM, while median floc size also decreases modestly from 42.0 to 38.8~}%DIFDELCMD < \textmu %%%
\DIFdel{m. The decline }\DIFdelend \DIFaddbegin \DIFadd{m$^{-3}$ is the organic matter density, $d_p$ is the primary particle diameter, and $\delta = 1~\mu$m is the assumed organic coating thickness.
}

\DIFadd{A key property of the }\citeA{nghiem2022mechanistic} \DIFadd{framework is that the predicted equilibrium floc diameter does not depend on the fractal dimension. In their model (Eq.~9 in their study), the fractal dimension $n_f$ enters both the aggregation and breakup terms of the floc growth equation through identical powers of $d_f$: aggregation scales as $d_f^{\,n_f+2}$ and breakup as $d_f^{\,n_f}$. When these two terms are set equal at equilibrium, $n_f$ cancels algebraically and the resulting equilibrium diameter $d_{f,\text{eq}}$ depends only on $\eta$ and $\theta$. Physically, this cancellation reflects the fact that the same fractal structure governs both how efficiently particles collide (aggregation) and how susceptible the aggregate is to turbulent stresses (breakup), so the two effects offset exactly at equilibrium.
}

\DIFadd{Using the calibrated exponents from }\citeA{nghiem2022mechanistic} \DIFadd{(their Eqs.~23--24), the equilibrium diameter scales as $d_{f,\text{eq}} \propto \eta^{r}[\theta^2(1-\theta)^2]^{q}$, where $r$ and $q$ are fitted constants. Since $\eta \propto u_*^{-3/4}$, the predicted shear dependence is $d_{f,\text{eq}} \propto u_*^{-0.75r}$, and the organic-matter dependence enters through $\theta$ (Eq.~\eqref{eq:theta}). Because the model parameters were calibrated for river field conditions rather than our tank geometry, we compare trends rather than absolute magnitudes: we normalize the predicted $d_{f,\text{eq}}$ by a single least-squares constant and plot it against the measured medians in Fig.~\ref{fig:6}a,b; the model captures the monotonic decrease of $d_f$ with shear and the weak non-monotonic response to organic matter.
}

\DIFadd{Although the equilibrium floc diameter is independent of $n_f$, the settling velocity is not. Equation~\eqref{eq:floc_settling} shows that $w_s$ depends on $n_f$ through two exponents: $d_p^{\,3-n_f}$ in the prefactor and $d_f^{\,n_f-1}$ in the size scaling. The }\citeA{nghiem2022mechanistic} \DIFadd{model adopts $n_f = 2$ throughout, which simplifies Eq.~\eqref{eq:floc_settling} to $w_s \propto d_p\, d_f$, a linear dependence on floc diameter in which $d_p$ and $b_1$ are absorbed into the prefactor. We refer to this as the base model. We compare the predicted and observed trends rather than absolute magnitudes for two reasons. First, the }\citeA{nghiem2022mechanistic} \DIFadd{model parameters were calibrated against river field data; applying them to our tank geometry, which has a different energy dissipation distribution, would require re-calibration. Second, the settling velocity prefactor contains the shape factor $b_1$, which is not independently constrained and may itself vary with $n_f$ (see Section~\ref{sec:dp}). We therefore normalize each model prediction by a single least-squares constant ($\hat{w}_{s,\text{base}} = k_{\text{base}}\, d_{f,\text{eq}}$) and compare the predicted trend against the measured medians in Fig.~\ref{fig:6}c,d. Because the base model is linear in $d_f$, and $d_f$ decreases monotonically with shear, the base model strictly predicts a monotonic decrease }\DIFaddend in $w_s$ with increasing \DIFdelbegin \DIFdel{OM is consistent with flocs becoming more porous and less dense (lower $n_f$)as organic binding promotes open, voluminous structures. Thus, both size and structure change with OM , and both contribute }\DIFdelend \DIFaddbegin \DIFadd{$u_*$. It therefore cannot reproduce the non-monotonic peak observed at $u_* = 0.03$~m\,s$^{-1}$ (Fig.~\ref{fig:6}c). Similarly, because $d_f$ increases slightly with OM (for $\theta < 0.5$), the base model predicts $w_s$ increasing with OM, opposite }\DIFaddend to the observed \DIFdelbegin \DIFdel{settling-velocity trends}\DIFdelend \DIFaddbegin \DIFadd{monotonic decline (Fig.~\ref{fig:6}d)}\DIFaddend .

These \DIFdelbegin \DIFdel{bulk trends motivate the stratified analysis that follows: the settling velocity depends not only on floc size but also on internal structure, and this }\DIFdelend \DIFaddbegin \DIFadd{discrepancies motivate a closer examination of the internal }\DIFaddend structural \DIFdelbegin \DIFdel{variability must be resolved to accurately characterize }\DIFdelend \DIFaddbegin \DIFadd{properties---the fractal dimension }\DIFaddend $n_f$ \DIFaddbegin \DIFadd{and the effective primary particle diameter $d_p$---that the base model holds fixed. In the following two subsections, we develop methods to extract condition-specific estimates of $n_f$ (Section~\ref{sec:stratifying}) and $d_p$ (Section~\ref{sec:dp}), and in Section~\ref{sec:reconciling} we show that incorporating this measured structural variability into the settling velocity prediction resolves the failures identified above}\DIFaddend .


\subsection{\DIFdelbegin \DIFdel{Fractal Dimension from Stratified and Aggregated Analyses}\DIFdelend \DIFaddbegin \DIFadd{Stratifying 2D Box-Counting Data}\DIFaddend }
\DIFdelbegin %DIFDELCMD < 

%DIFDELCMD < %%%
\subsubsection{\DIFdel{Stratification Methodology and Bin Selection}}
%DIFAUXCMD
\addtocounter{subsubsection}{-1}%DIFAUXCMD
\DIFdelend \DIFaddbegin \label{sec:stratifying}
\DIFaddend 

To infer a \DIFdelbegin \DIFdel{three-dimensional }\DIFdelend fractal dimension from settling data, we group flocs into bins based on their \DIFdelbegin \DIFdel{two-dimensional }\DIFdelend \DIFaddbegin \DIFadd{2D }\DIFaddend box-counting \DIFdelbegin \DIFdel{fractal dimension, $n_f^{(2\mathrm{D})}$}\DIFdelend \DIFaddbegin \DIFadd{dimension, $n_f^{(\mathrm{2D})}$}\DIFaddend , measured from images. The bin width in \DIFdelbegin \DIFdel{$n_f^{(2\mathrm{D})}$ is }\DIFdelend \DIFaddbegin \DIFadd{$n_f^{(\mathrm{2D})}$ serves only as }\DIFaddend a coarse-graining scale \DIFdelbegin \DIFdel{rather than }\DIFdelend \DIFaddbegin \DIFadd{for separating discrete structural classes, not as }\DIFaddend a physical control parameter. If a population contains \DIFdelbegin \DIFdel{discrete }\DIFdelend \DIFaddbegin \DIFadd{distinct }\DIFaddend structural classes (e.g., non-flocculating grains with $n_f \approx 3$ and flocs with $n_f \approx 2$), finer stratification \DIFdelbegin \DIFdel{simply produces }\DIFdelend \DIFaddbegin \DIFadd{merely subdivides these classes into }\DIFaddend redundant bins that recover the same class-specific slopes. 
\DIFdelbegin \DIFdel{In finite datasets, the practical objective is therefore to resolve statistically distinct structural subpopulations while retaining enough particles per bin to constrain the conditional $w_s$--$d_f$ slope.
}\DIFdelend 

As \DIFdelbegin \DIFdel{an initial test case for flocs using our stratification methodology}\DIFdelend \DIFaddbegin \DIFadd{a sensitivity analysis}\DIFaddend , we computed the \DIFdelbegin \DIFdel{two-dimensional }\DIFdelend \DIFaddbegin \DIFadd{2D }\DIFaddend box-counting \DIFdelbegin \DIFdel{fractal dimension, $n_f^{(2\mathrm{D})}$}\DIFdelend \DIFaddbegin \DIFadd{dimension, $n_f^{(\mathrm{2D})}$}\DIFaddend , for Experiment~3 (2~wt\% organic matter; shear velocity $u_* = 0.02~\mathrm{m\,s^{-1}}$). To evaluate the impact of stratification, we divided the dataset into bins of width 0.1, 0.05, and 0.025, corresponding to 3, 6, and 12 bins of \DIFdelbegin \DIFdel{$n_f^{(2\mathrm{D})}$}\DIFdelend \DIFaddbegin \DIFadd{$n_f^{(\mathrm{2D})}$}\DIFaddend , respectively, within a manually defined range from 1.3 to 1.6. Within each bin, we fitted a power-law relation between settling velocity, $w_s$, and floc diameter, $d_f$, in log--log space (Figs.~\DIFdelbegin \DIFdel{\ref{fig:6}}\DIFdelend \DIFaddbegin \DIFadd{\ref{fig:7}}\DIFaddend a, d, g). The model $\log_{10}(w_s) = m \log_{10}(d_f) + C$ was calibrated using an ensemble Markov chain Monte Carlo (MCMC) sampler to obtain robust parameter estimates and verify convergence. After burn-in and thinning, the posterior medians of the slope $m$ were converted to the hydrodynamically relevant \DIFdelbegin \DIFdel{three-dimensional }\DIFdelend \DIFaddbegin \DIFadd{3D }\DIFaddend fractal dimension, denoted simply as $n_f$, using the relation $n_f = m + 1$.

Plotting the inferred \DIFdelbegin \DIFdel{three-dimensional }\DIFdelend \DIFaddbegin \DIFadd{3D }\DIFaddend fractal dimension, $n_f$, against the corresponding group-mean \DIFdelbegin \DIFdel{two-dimensional value, $n_f^{(2\mathrm{D})}$}\DIFdelend \DIFaddbegin \DIFadd{2D value, $n_f^{(\mathrm{2D})}$}\DIFaddend , shows how the stratification behaves across binning schemes. A linear trend captures the binned data reasonably well for all cases (Figs.~\DIFdelbegin \DIFdel{\ref{fig:6}}\DIFdelend \DIFaddbegin \DIFadd{\ref{fig:7}}\DIFaddend c, f, i). However, because the number of bins is small (3, 6, or 12), the resulting coefficients of determination are not themselves a robust metric of model adequacy, as high $R^2$ values are expected even in the presence of substantial uncertainty.

A more informative assessment comes from quantifying the dynamical consequences of refining the stratification. From Eq.~\eqref{eq:floc_settling}, assuming the log--log regression lines converge near the primary particle diameter $d_p$, a change $\Delta n_f$ produces a vertical shift $\Delta \log_{10}(w_s) \approx \Delta n_f \log_{10}(d_f/d_p)$. Over the observed range of macroscopic floc sizes ($d_f/d_p \approx 10^1$--$10^2$), a change of $\Delta n_f = 0.01$ produces a shift of $\Delta \log_{10}(w_s) \approx 0.01$--$0.02$, corresponding to a change in settling velocity of less than 5\%. This effect is minimal relative to measurement scatter. Consequently, increasingly fine stratification does not reveal qualitatively different settling regimes but instead primarily inflates variance in the estimated slopes as bin populations decrease. We therefore adopt a bin width of 0.1 for subsequent analyses. This choice resolves meaningful structural heterogeneity while maintaining sufficient sample sizes for regression, and avoids over-resolving differences that have negligible impact on settling velocities over the observed floc-size range.

\begin{figure}[htbp]
    \centering
    \DIFdelbeginFL \DIFdelFL{\includegraphics[width=1\textwidth]{figures/6.pdf}
    }\DIFdelendFL \DIFaddbeginFL \DIFaddFL{\includegraphics[width=1\textwidth]{figures/fig7.pdf}
    }\DIFaddendFL \caption{Stratification of floc data by \DIFdelbeginFL \DIFdelFL{two-dimensional (}\DIFdelendFL 2D \DIFdelbeginFL \DIFdelFL{) }\DIFdelendFL box-counting \DIFdelbeginFL \DIFdelFL{fractal }\DIFdelendFL dimension (\DIFdelbeginFL \DIFdelFL{\(n_f^{(2\mathrm{D})}\)}\DIFdelendFL \DIFaddbeginFL \DIFaddFL{\(n_f^{(\mathrm{2D})}\)}\DIFaddendFL ) for Experiment~3, illustrating how binning strategy influences the estimation of \DIFdelbeginFL \DIFdelFL{three-dimensional }\DIFdelendFL \DIFaddbeginFL \DIFaddFL{3D }\DIFaddendFL fractal dimensions. (a--c) Analysis using 3 bins: (a) log--log plots of settling velocity (\(w_s\)) versus floc diameter (\(d_f\)), color-coded by 2D \DIFdelbeginFL \DIFdelFL{fractal-dimension }\DIFdelendFL \DIFaddbeginFL \DIFaddFL{box-counting dimension }\DIFaddendFL bin; (b) fitted slopes (\(m\)) for each bin; (c) resulting \DIFdelbeginFL \DIFdelFL{three-dimensional }\DIFdelendFL \DIFaddbeginFL \DIFaddFL{3D }\DIFaddendFL fractal dimension (\(n_f = m + 1\)) plotted against mean \DIFdelbeginFL \DIFdelFL{\(n_f^{(2\mathrm{D})}\) }\DIFdelendFL \DIFaddbeginFL \DIFaddFL{\(n_f^{(\mathrm{2D})}\) }\DIFaddendFL for each bin, showing a strong linear relationship. (d--f) Same analysis using 6 bins; (g--i) using 12 bins, highlighting increased scatter and reduced regression stability with finer binning.}
    \DIFdelbeginFL %DIFDELCMD < \label{fig:6}
%DIFDELCMD < %%%
\DIFdelendFL \DIFaddbeginFL \label{fig:7}
\DIFaddendFL \end{figure}

\DIFdelbegin \subsubsection{\DIFdel{3D Fractal Dimension Across All Experiments}}
%DIFAUXCMD
\addtocounter{subsubsection}{-1}%DIFAUXCMD
%DIFDELCMD < 

%DIFDELCMD < %%%
\DIFdel{We expanded the analysis to all floc datasets (Experiments ~}\DIFdelend \DIFaddbegin \DIFadd{The six floc experiments (Experiments }\DIFaddend 3--8) \DIFaddbegin \DIFadd{comprise approximately 12}{\DIFadd{,}}\DIFadd{500 individual floc particles}\DIFaddend . Within each 0.1-wide \DIFdelbegin \DIFdel{\(n_f^{(2\mathrm{D})}\) bin, anomalous measurements were objectively identified and removed using an interquartile range (IQR) filter applied to the settling velocities: any particle with $\log_{10}(w_s)$ outside \([Q_1 - 1.5\,\mathrm{IQR},\ Q_3 + 1.5\,\mathrm{IQR}]\), where \(\mathrm{IQR} = Q_3 - Q_1\), was excluded. Only bins retaining at least 50 particles after filtering were analyzed, yielding a typical count of five to nine }\DIFdelend \DIFaddbegin \DIFadd{$n_f^{(\mathrm{2D})}$ bin, outliers were removed using the interquartile range criterion }[\DIFadd{$Q_1 - 1.5\,\mathrm{IQR},\, Q_3 + 1.5\,\mathrm{IQR}$}]\DIFadd{, yielding five to eight }\DIFaddend bins per experiment \DIFdelbegin \DIFdel{. Across experiments, analyzed bins typically contained tens to hundreds of flocs (median $\approx 280$ across bins with $n\ge50$), with peak counts concentrated near intermediate fractal dimensions (\(n_f^{(2\mathrm{D})}\approx1.4\!-\!1.6\)); bins near the distribution tails (\(n_f^{(2\mathrm{D})}\lesssim1.2\) or \(\gtrsim1.8\)) contained substantially fewer particles and were excluded by the minimum-count criterion. For each bin, the posterior median slope \(m\) from the log--log regression was converted to the corresponding three-dimensional fractal dimension as \(n_f = m + 1\). Inset panels in Fig.
~\ref{fig:7} show the derived 3D \(n_f\) versus the bin-averaged \(n_f^{(2\mathrm{D})}\). Across all six experiments, stratified results reveal a strong positive correlation between these metrics, consistently supported by high \(R^2\) values.
}\DIFdelend \DIFaddbegin \DIFadd{with a median population of 190 flocs per bin. The same MCMC regression procedure described above was then applied within each bin to obtain $n_f$.
}\DIFaddend 

\DIFdelbegin \DIFdel{The stratified analysis yields experiment-level median $n_f$ values that vary systematically with forcing. For the shear series (Fig.~\ref{fig:7}}\DIFdelend \DIFaddbegin \DIFadd{Figure~\ref{fig:8} shows the log--log relationship between settling velocity $w_s$ and floc diameter $d_f$ for experiments with increasing shear velocity (Figs.~\ref{fig:8}}\DIFaddend a--c) \DIFdelbegin \DIFdel{, median $n_f$ increases from 1.77 at $u_* = 0.02$~m\,s$^{-1}$ to 1.95 at $u_* = 0.03$~m\,s$^{-1}$ and 2.12 at $u_* = 0.04$~m\,s$^{-1}$ (stratified IQRs of 0.26, 0.08, and 0.10, respectively). For the organic-matter series (Fig.~\ref{fig:7}}\DIFdelend \DIFaddbegin \DIFadd{and increasing organic matter content (Figs.~\ref{fig:8}}\DIFaddend d--f)\DIFdelbegin \DIFdel{, median }\DIFdelend \DIFaddbegin \DIFadd{. Within each panel, power-law regressions fitted to individual $n_f^{(\mathrm{2D})}$ bins reveal distinct slopes that vary systematically with floc structure, and the corresponding 3D fractal dimensions }\DIFaddend $n_f$ \DIFdelbegin \DIFdel{decreases from 1.84 at 1~wt\% OM to 1.66 at 2~wt\% OM and 1.63 at 3~wt\% OM (stratified IQRs of 0.14, 0.15, and 0.18). In contrast, the aggregated (bulk) analysis yields comparatively invariant median }\DIFdelend \DIFaddbegin \DIFadd{are obtained via the relation $n_f = m + 1$. As the $n_f^{(\mathrm{2D})}$ bin midpoint increases, the recovered }\DIFaddend $n_f$ \DIFdelbegin \DIFdel{values: 1.98, 1.94, and 2.00 for the shear series, and 1.80, 1.74, and 1.75 for the OM series (all with narrow posterior IQRs of 0.03--0.04). The aggregated IQRs reflect only the posterior uncertainty of a single bulk $w_s$--$d_f$ regression , which concentrates with many observations; consequently}\DIFdelend \DIFaddbegin \DIFadd{also increases monotonically, demonstrating that the 2D structural signature can be used to resolve distinct 3D fractal dimensions within a mixed population. The inset panels illustrate this trend, plotting the bin-median $n_f$ against the $n_f^{(\mathrm{2D})}$ bin midpoint with 95\% credible intervals derived from the MCMC posteriors. For comparison, the aggregated (unstratified) regression and its 95\% credible interval are shown as the orange dashed line and shaded band, respectively. Notably}\DIFaddend , the aggregated \DIFdelbegin \DIFdel{approach assigns negligible probability to the lower }\DIFdelend \DIFaddbegin \DIFadd{credible interval does not capture the range of }\DIFaddend $n_f$ values \DIFdelbegin \DIFdel{resolved }\DIFdelend \DIFaddbegin \DIFadd{captured }\DIFaddend by the stratified \DIFdelbegin \DIFdel{method}\DIFdelend \DIFaddbegin \DIFadd{bins, indicating that a single power-law regression applied to the bulk population cannot capture the structural heterogeneity present within each experiment}\DIFaddend .

\begin{figure}[htbp]
    \centering
    \DIFdelbeginFL \DIFdelFL{\includegraphics[width=1\textwidth]{figures/flocfigure7.pdf}
    }\DIFdelendFL \DIFaddbeginFL \DIFaddFL{\includegraphics[width=1\textwidth]{figures/fig8.pdf}
    }\DIFaddendFL \caption{Comparison of stratified and aggregated (bulk) fits for all floc datasets. (a--c) Varying shear velocities of 0.02, 0.03, and 0.04~m\,s$^{-1}$. (d--f) Varying organic matter contents of 1\%, 2\%, and 3\%. Inset plots show the relationship between \DIFdelbeginFL \DIFdelFL{\(n_f^{(2\mathrm{D})}\) }\DIFdelendFL \DIFaddbeginFL \DIFaddFL{\(n_f^{(\mathrm{2D})}\) }\DIFaddendFL and slope-inferred \(n_f\) for each experiment; points represent bin medians (with 95\% credible intervals) and are color-coded by \DIFdelbeginFL \DIFdelFL{\(n_f^{(2\mathrm{D})}\) }\DIFdelendFL \DIFaddbeginFL \DIFaddFL{\(n_f^{(\mathrm{2D})}\) }\DIFaddendFL group to make within-population variability visually explicit.}
    \DIFdelbeginFL %DIFDELCMD < \label{fig:7}
%DIFDELCMD < %%%
\DIFdelendFL \DIFaddbeginFL \label{fig:8}
\DIFaddendFL \end{figure}


\DIFdelbegin \subsubsection{\DIFdel{Effects of Shear Velocity and Organic Matter}}
%DIFAUXCMD
\addtocounter{subsubsection}{-1}%DIFAUXCMD
%DIFDELCMD < 

%DIFDELCMD < %%%
\DIFdelend The dependence of the stratified and aggregated $n_f$ estimates on \DIFdelbegin \DIFdel{environmental forcing }\DIFdelend \DIFaddbegin \DIFadd{shear velocity and organic matter }\DIFaddend is summarized in Fig.~\DIFdelbegin \DIFdel{\ref{fig:8}}\DIFdelend \DIFaddbegin \DIFadd{\ref{fig:9}}\DIFaddend a,b. \DIFaddbegin \DIFadd{In the stratified analysis, the interquartile range (IQR) is computed directly from the population of inferred $n_f$ values across bins, weighted by the number of particles contributing to each bin, because each bin represents an independent subset of flocs with distinct fractal structure. In contrast, in the aggregated analysis the IQR is derived from the Bayesian posterior distribution of the fitted slope parameter, since all particles are pooled into a single regression and variability arises from parameter uncertainty rather than particle-to-particle heterogeneity.
}

\DIFaddend Stratified median $n_f$ increases approximately linearly with shear velocity at a rate of $\Delta n_f / \Delta u_* \approx 17$~s\,m$^{-1}$ ($n_f = 17.4\,u_* + 1.42$\DIFdelbegin \DIFdel{; }\DIFdelend \DIFaddbegin \DIFadd{, }\DIFaddend Fig.~\DIFdelbegin \DIFdel{\ref{fig:8}}\DIFdelend \DIFaddbegin \DIFadd{\ref{fig:9}}\DIFaddend a), consistent with shear-driven breakup preferentially destroying loose, low-$n_f$ flocs and favoring compact, high-$n_f$ \DIFdelbegin \DIFdel{survivors}\DIFdelend \DIFaddbegin \DIFadd{flocs}\DIFaddend . The aggregated analysis yields a nearly flat response ($n_f = 1.0\,u_* + 1.94$), demonstrating that the shear-driven structural trend is effectively invisible in bulk data. \DIFdelbegin %DIFDELCMD < 

%DIFDELCMD < %%%
\DIFdelend For organic matter (Fig.~\DIFdelbegin \DIFdel{\ref{fig:8}}\DIFdelend \DIFaddbegin \DIFadd{\ref{fig:9}}\DIFaddend b), stratified median $n_f$ decreases at $\Delta n_f / \Delta \mathrm{OM} \approx -0.11$ per wt\% ($n_f = -0.11\,\mathrm{OM} + 1.92$), indicating that organic binding promotes open, porous structures. The aggregated trend is \DIFaddbegin \DIFadd{again }\DIFaddend much weaker ($n_f = -0.03\,\mathrm{OM} + 1.81$). \DIFdelbegin \DIFdel{Stratified IQRs also widen with increasing OM (from 0.14 to 0.18), indicating that organic-rich populations contain enhanced structural diversity}\DIFdelend \DIFaddbegin \DIFadd{This insensitivity arises because the aggregated regression fits a single power law through a heterogeneous mixture of floc structures, each with a different $w_s$--$d_f$ slope. The resulting best-fit slope is dominated by the most populated $n_f^{(\mathrm{2D})}$ bins and by the large scatter inherent in pooling structurally distinct subpopulations, which masks the systematic variation that the stratified analysis resolves. In effect, the noise introduced by collapsing all fractal classes into one regression absorbs the signal that distinguishes compact from porous flocs}\DIFaddend .

\begin{figure}[htbp]
    \centering
    \DIFdelbeginFL \DIFdelFL{\includegraphics[width=0.8\textwidth]{figures/flocpaperfigure8.pdf}
    }\DIFdelendFL \DIFaddbeginFL \DIFaddFL{\includegraphics[width=1\textwidth]{figures/fig9.pdf}
    }\DIFaddendFL \caption{Comparison of stratified and aggregated (bulk) analyses of floc fractal dimensions. (a) Median inferred \(n_f\) versus shear velocity. (b) Median \(n_f\) versus organic matter (OM) content. \DIFdelbeginFL \DIFdelFL{Boxes }\DIFdelendFL \DIFaddbeginFL \DIFaddFL{Error bars }\DIFaddendFL denote interquartile ranges (IQRs) and \DIFdelbeginFL \DIFdelFL{horizontal ticks }\DIFdelendFL \DIFaddbeginFL \DIFaddFL{markers }\DIFaddendFL indicate medians; for stratified estimates \DIFaddbeginFL \DIFaddFL{(orange) }\DIFaddendFL the IQR reflects particle-scale variability across bins weighted by population, whereas for aggregated estimates \DIFaddbeginFL \DIFaddFL{(navy) }\DIFaddendFL it reflects uncertainty in the Bayesian posterior of a single fitted parameter. (c) \DIFdelbeginFL \DIFdelFL{Relationship between three-dimensional (\(n_f\)) and two-dimensional (\(n_f^{(2\mathrm{D})}\)) box-counting fractal dimensions across all conditions, colored by experiment type. (d) }\DIFdelendFL Mapping between 2D and 3D fractal dimensions: \DIFdelbeginFL \DIFdelFL{blue }\DIFdelendFL \DIFaddbeginFL \DIFaddFL{black }\DIFaddendFL dashed line, theoretical relation \DIFdelbeginFL \DIFdelFL{\(n_f = n_f^{(2\mathrm{D})} + 1\)}\DIFdelendFL \DIFaddbeginFL \DIFaddFL{\(n_f = n_f^{(\mathrm{2D})} + 1\)}\DIFaddendFL ; \DIFdelbeginFL \DIFdelFL{red }\DIFdelendFL \DIFaddbeginFL \DIFaddFL{orange }\DIFaddendFL solid line, empirical 2D box-counting to 3D box-counting (BC) conversion (\DIFdelbeginFL \DIFdelFL{$n_f = 0.81\,(n_f^{(2\mathrm{D})})^{1.81}$}\DIFdelendFL \DIFaddbeginFL \DIFaddFL{$n_f = 0.81\,(n_f^{(\mathrm{2D})})^{1.81}$}\DIFaddendFL ); \DIFdelbeginFL \DIFdelFL{red }\DIFdelendFL \DIFaddbeginFL \DIFaddFL{orange }\DIFaddendFL dashed line, empirical 2D box-counting to 3D power-law (PL) conversion (\DIFdelbeginFL \DIFdelFL{$n_f = 0.20\,(n_f^{(2\mathrm{D})})^{4.08}$}\DIFdelendFL \DIFaddbeginFL \DIFaddFL{$n_f = 0.20\,(n_f^{(\mathrm{2D})})^{4.08}$}\DIFaddendFL ) \cite{wang2022fractal}. Experimental data (\DIFdelbeginFL \DIFdelFL{black }\DIFdelendFL \DIFaddbeginFL \DIFaddFL{grey }\DIFaddendFL markers) align most closely with the empirical box-counting model.}
    \DIFdelbeginFL %DIFDELCMD < \label{fig:8}
%DIFDELCMD < %%%
\DIFdelendFL \DIFaddbeginFL \label{fig:9}
\DIFaddendFL \end{figure}

\DIFdelbegin \subsubsection{\DIFdel{Relationship Between 2D and 3D Fractal Dimensions}}
%DIFAUXCMD
\addtocounter{subsubsection}{-1}%DIFAUXCMD
%DIFDELCMD < 

%DIFDELCMD < %%%
\DIFdel{Our experimental results (Fig.~\ref{fig:8}c ) reveal a consistent positive relationship between the two- }\DIFdelend \DIFaddbegin \DIFadd{Figure~\ref{fig:9}c collects the per-bin data from all six experiments to examine the mapping between $n_f^{(\mathrm{2D})}$ }\DIFaddend and \DIFdelbegin \DIFdel{three-dimensional fractal dimensions across all conditions. Plotting the stratified median \(n_f\) against the binned \(n_f^{(2\mathrm{D})}\) shows that denser, more compact structures (higher \(n_f^{(2\mathrm{D})}\)) systematically exhibit larger \(n_f\). This trend holds regardless of whether the grouping parameter is shear velocity or organic matter (OM) content.
}%DIFDELCMD < 

%DIFDELCMD < %%%
\DIFdel{To assess the underlying physical relationship, we compared our results to theoretical and empirical conversion models(Fig.~\ref{fig:8}d)}\DIFdelend \DIFaddbegin \DIFadd{$n_f$. The 40 bin-median values (grey circles) are compared against three reference models}\DIFaddend . The idealized \DIFdelbegin \DIFdel{linear relationship for perfectly space-filling Euclidean objects, \(n_f = n_f^{(2\mathrm{D})} + 1\) (blue dashedline), correctly describes smooth solids such as a sphere (\(n_f^{(2\mathrm{D})}=2\), \(n_f=3\)) but overpredicts the structural dimensions of porous flocs}\DIFdelend \DIFaddbegin \DIFadd{Euclidean relation $n_f = n_f^{(\mathrm{2D})} + 1$ (black dashed) assumes perfect space-filling and systematically overpredicts $n_f$ across the measured range}\DIFaddend . The empirical \DIFdelbegin \DIFdel{models of }%DIFDELCMD < \citeA{wang2022fractal}%%%
\DIFdel{---a }\DIFdelend \DIFaddbegin \DIFadd{power-law (PL) and }\DIFaddend box-counting (BC) \DIFdelbegin \DIFdel{conversion, \(n_f = 0.81\,(n_f^{(2\mathrm{D})})^{1.81}\) (red solid line) , and a power-law (PL ) conversion, \(n_f = 0.20\,(n_f^{(2\mathrm{D})})^{4.08}\) (red dashed line) ---provide better descriptions}\DIFdelend \DIFaddbegin \DIFadd{conversions of }\citeA{wang2022fractal}\DIFadd{, calibrated on computationally generated aggregates, bracket the data more closely. The BC model ($n_f = 0.8118\, {n_f^{(\mathrm{2D})}}^{1.8054}$, solid orange) tracks the upper envelope of the observations, while the PL model ($n_f = 0.2015\, {n_f^{(\mathrm{2D})}}^{4.079}$, dashed orange) underpredicts at moderate $n_f^{(\mathrm{2D})}$}\DIFaddend . Our data \DIFdelbegin \DIFdel{align most closely with the BC conversion, }\DIFdelend \DIFaddbegin \DIFadd{scatter between these two empirical curves, with the closest overall agreement to the BC conversion. This alignment is }\DIFaddend consistent with the \DIFdelbegin \DIFdel{sub-Euclidean space filling expected of porous aggregates: a }\DIFdelend \DIFaddbegin \DIFadd{fact that both our stratification variable and the }\citeA{wang2022fractal} \DIFadd{BC model are based on box-counting dimensions. Mismatch may arise because the }\citeA{wang2022fractal} \DIFadd{model constructs flocs from monodisperse primary particles of uniform size and shape, whereas natural flocs comprise polydisperse, irregularly shaped constituents whose heterogeneous packing alters the }\DIFaddend 2D \DIFdelbegin \DIFdel{projection of a fractal object loses structural information, so the 3D dimension grows sublinearly rather than by a constant offset of unity. Although the data appear approximately linear within our observed range of \(n_f^{(2\mathrm{D})}\) (Fig.~\ref{fig:8}c), the underlying relationship is inherently nonlinear, as confirmed by the better agreement with the power-law BC conversion than with the linear Euclidean model}\DIFdelend \DIFaddbegin \DIFadd{projection geometry}\DIFaddend .

\subsection{Inferring Primary Particle Diameter ($d_p$)}
\DIFaddbegin \label{sec:dp}
\DIFaddend 

\DIFdelbegin \DIFdel{To estimate the primary particle diameter ($d_p$), we used a Bayesian approach to quantify uncertainty in the intersection method. This approach assumes that }\DIFdelend \DIFaddbegin \DIFadd{The stratified regressions also provide a means of estimating }\DIFaddend the effective primary particle \DIFdelbegin \DIFdel{size controlling the settling relation is shared across the $n_f$-stratified subpopulations within a given run; accordingly, the inferred }\DIFdelend \DIFaddbegin \DIFadd{diameter, }\DIFaddend $d_p$\DIFdelbegin \DIFdel{should be interpreted as an effective parameter in }\DIFdelend \DIFaddbegin \DIFadd{. In log--log space, }\DIFaddend Eq.~\eqref{eq:floc_settling} \DIFdelbegin \DIFdel{rather than a direct census of the finest grains present in the source material. For each stratified floc group, we drew samples from the posterior distribution of the log--log settling fits and computed }\DIFdelend \DIFaddbegin \DIFadd{predicts that regression lines for different $n_f$ converge as $d_f \to d_p$: at the primary-particle scale, a floc reduces to a single grain whose settling velocity is independent of fractal structure. We exploit this convergence by computing }\DIFaddend all pairwise intersections \DIFdelbegin \DIFdel{between fitted lines in log--log space. Under the constant-$b_1$ assumption (using $b_1 = 18$, the theoretical Stokes value for smooth spheres, as a reference baseline), the resulting intersection diameters $d_\text{intersect}$ provide an estimate of an effective $d_p$ (Eq.~\eqref{eq:dpintersectionideal}); allowing }\DIFdelend \DIFaddbegin \DIFadd{of the per-bin regression lines from their MCMC posteriors and assume a constant shape factor }\DIFaddend $b_1$\DIFdelbegin \DIFdel{to vary among groups yields the more general relation in Eq.~\eqref{eq:dpintersection}. Repeating this calculation across all }\DIFdelend \DIFaddbegin \DIFadd{. Repeating the calculation across thousands of }\DIFaddend posterior draws yields an ensemble of intersection diameters \DIFdelbegin \DIFdel{; the resulting posterior distribution is what we plot as the inferred $d_p$ distribution in Fig.~\ref{fig:9}. All regression lines used in this intersection analysis satisfy the minimum sample-size criterion ($n\ge50$ flocs per group), and the posterior spread therefore }\DIFdelend \DIFaddbegin \DIFadd{whose spread }\DIFaddend reflects both measurement scatter and finite-sample uncertainty. \DIFdelbegin %DIFDELCMD < 

%DIFDELCMD < %%%
\DIFdel{The intersection analysis yields inferred }\DIFdelend \DIFaddbegin \DIFadd{Fig.~\ref{fig:10}a--c compare the resulting }\DIFaddend $d_p$ \DIFdelbegin \DIFdel{values of 21.7, 23.0, and 29.8~}%DIFDELCMD < \textmu %%%
\DIFdel{m for the low-, mid-, and high-shear }\DIFdelend \DIFaddbegin \DIFadd{posterior (orange) with the known input grain-size distribution (gray) for the three shear-velocity }\DIFaddend experiments, and \DIFdelbegin \DIFdel{34.2, 27.9, and 17.6~}%DIFDELCMD < \textmu %%%
\DIFdel{m for the 1\%, 2\%, and 3\% OM }\DIFdelend \DIFaddbegin \DIFadd{Fig.~\ref{fig:10}e--g show the corresponding comparison for the three organic-matter }\DIFaddend experiments. 

\begin{figure}[htbp]
    \centering
    \DIFdelbeginFL \DIFdelFL{\includegraphics[width=1\textwidth]{figures/9.pdf}
    }\DIFdelendFL \DIFaddbeginFL \DIFaddFL{\includegraphics[width=1\textwidth]{figures/fig10.pdf}
    }\DIFaddendFL \caption{
        \DIFaddbeginFL \DIFaddFL{Inferred primary particle diameter ($d_p$) distributions and summary statistics. (a--c) }\DIFaddendFL Cumulative percent finer by volume \DIFdelbeginFL \DIFdelFL{distributions }\DIFdelendFL comparing \DIFaddbeginFL \DIFaddFL{the }\DIFaddendFL known input \DIFdelbeginFL \DIFdelFL{and }\DIFdelendFL \DIFaddbeginFL \DIFaddFL{grain-size distribution (grey) with the }\DIFaddendFL inferred \DIFaddbeginFL \DIFaddFL{$d_p$ }\DIFaddendFL posterior \DIFdelbeginFL \DIFdelFL{effective primary particle diameters }\DIFdelendFL (\DIFdelbeginFL \DIFdelFL{$d_p$}\DIFdelendFL \DIFaddbeginFL \DIFaddFL{orange}\DIFaddendFL ) for \DIFdelbeginFL \DIFdelFL{all experiments}\DIFdelendFL \DIFaddbeginFL \DIFaddFL{each shear-velocity experiment; KS statistics quantify agreement}\DIFaddendFL . \DIFdelbeginFL \DIFdelFL{Panels show }\DIFdelendFL (\DIFdelbeginFL \DIFdelFL{a--c}\DIFdelendFL \DIFaddbeginFL \DIFaddFL{d}\DIFaddendFL ) \DIFdelbeginFL \DIFdelFL{varying }\DIFdelendFL \DIFaddbeginFL \DIFaddFL{Median inferred $d_p$ (orange, with IQR error bars) versus }\DIFaddendFL shear velocity\DIFaddbeginFL \DIFaddFL{, alongside the input-sediment median }\DIFaddendFL and \DIFaddbeginFL \DIFaddFL{IQR }\DIFaddendFL (\DIFdelbeginFL \DIFdelFL{d--f}\DIFdelendFL \DIFaddbeginFL \DIFaddFL{grey}\DIFaddendFL )\DIFdelbeginFL \DIFdelFL{varying OM content}\DIFdelendFL . \DIFdelbeginFL \DIFdelFL{The Kolmogorov--Smirnov }\DIFdelendFL (\DIFdelbeginFL \DIFdelFL{KS}\DIFdelendFL \DIFaddbeginFL \DIFaddFL{e--g}\DIFaddendFL ) \DIFdelbeginFL \DIFdelFL{statistic quantifies the agreement between }\DIFdelendFL \DIFaddbeginFL \DIFaddFL{Same as (a--c) for }\DIFaddendFL the \DIFdelbeginFL \DIFdelFL{two distributions in each panel}\DIFdelendFL \DIFaddbeginFL \DIFaddFL{three organic-matter experiments. (h) Median inferred $d_p$ versus organic matter content}\DIFaddendFL .
    }
    \DIFdelbeginFL %DIFDELCMD < \label{fig:9}
%DIFDELCMD < %%%
\DIFdelendFL \DIFaddbeginFL \label{fig:10}
\DIFaddendFL \end{figure}

\DIFdelbegin \DIFdel{Fig.~\ref{fig:9} compares the cumulative percent finer by volume distributions for the known input particle sizes and }\DIFdelend \DIFaddbegin \DIFadd{Results show agreement between }\DIFaddend the inferred $d_p$ \DIFdelbegin \DIFdel{from the intersection approach}\DIFdelend \DIFaddbegin \DIFadd{distribution from flocs and the input grain-size distribution, particularly in their central tendency (e.g., alignment of medians) (Fig.~\ref{fig:10})}\DIFaddend . The Kolmogorov--Smirnov (KS) statistic \DIFdelbegin \DIFdel{, reported in each panel, quantifies the maximum difference between the two cumulative distributions: }\DIFdelend \DIFaddbegin \DIFadd{ranges from 0.28 to 0.43 across all six experiments, indicating moderate agreement (}\DIFaddend a KS value of 0 \DIFdelbegin \DIFdel{indicates }\DIFdelend \DIFaddbegin \DIFadd{denotes }\DIFaddend perfect agreement, \DIFdelbegin \DIFdel{while }\DIFdelend \DIFaddbegin \DIFadd{whereas }\DIFaddend a value near 1 \DIFdelbegin \DIFdel{is a complete mismatch. Values between 0.29 and 0.41, as observed across all six experimental conditions, indicate moderate agreement , particularly in the central tendency (e.g., alignment of medians).
}%DIFDELCMD < 

%DIFDELCMD < %%%
\DIFdel{The input grain size }\DIFdelend \DIFaddbegin \DIFadd{indicates complete mismatch). Perfect agreement is not expected because the input grain-size }\DIFaddend distribution in Fig.~\DIFdelbegin \DIFdel{\ref{fig:9} shows }\DIFdelend \DIFaddbegin \DIFadd{\ref{fig:10} includes }\DIFaddend all particle sizes added to the experiments, but not all of these necessarily serve as primary particles in flocs. \DIFdelbegin \DIFdel{Some input size fractions---such as the finest clays or coarsest silts and sands---may not be equally incorporated into flocs }\DIFdelend \DIFaddbegin \DIFadd{In every case the inferred $d_p$ distribution is substantially narrower than the input grain distribution, with posteriors peaking in the coarse-silt range ($\sim$10--40~}\textmu \DIFadd{m), suggesting that the effective building blocks of the observed flocs are drawn predominantly from the silt fraction }\cite{nghiem2026flocculated}\DIFadd{.
}

\DIFadd{Panels (d) and (h) of Fig.~\ref{fig:10} summarize the median $d_p$ (orange, with IQR error bars) alongside the input-sediment median and IQR (grey) across shear velocity and organic matter. With increasing shear velocity (panel d), the inferred $d_p$ increases from 21.7~}\textmu \DIFadd{m at $u_* = 0.02$~m\,s$^{-1}$ to 23.0~}\textmu \DIFadd{m at 0.03~m\,s$^{-1}$ and 29.8~}\textmu \DIFadd{m at 0.04~m\,s$^{-1}$, consistent with stronger shear selectively destroying the most fragile aggregates and shifting the effective primary scale toward coarser grains. With increasing organic matter (panel h), the median $d_p$ decreases monotonically from 34.2~}\textmu \DIFadd{m at 1\% OM to 27.9~}\textmu \DIFadd{m at 2\% and 17.6~}\textmu \DIFadd{m at 3\%, indicating that organic binding enables finer particles to participate as effective building blocks. Across all conditions the inferred $d_p$ remains well below the input-sediment median of $\sim$39~}\textmu \DIFadd{m (grey markers in panels d and h)}\DIFaddend . 
\DIFdelbegin \DIFdel{Consequently, differences between the input grain size and inferred primary particle (}\DIFdelend \DIFaddbegin 

\DIFadd{The inferred }\DIFaddend $d_p$ \DIFdelbegin \DIFdel{) distributions may }\DIFdelend \DIFaddbegin \DIFadd{thus represents an effective primary scale associated with the population of flocs contributing to the measured settling velocities, rather than the full input grain-size distribution. Differences between the input grain-size and inferred $d_p$ distributions }\DIFaddend reflect both uncertainties in the intersection method and real selective incorporation of certain \DIFdelbegin \DIFdel{sizes }\DIFdelend \DIFaddbegin \DIFadd{size classes }\DIFaddend into flocs. Prolonged \DIFdelbegin \DIFdel{cloudiness after }\DIFdelend \DIFaddbegin \DIFadd{turbidity following }\DIFaddend the experiments indicates that a fraction of very fine material remained in suspension \DIFaddbegin \DIFadd{and was not flocculated}\DIFaddend ; however, this observation alone \DIFdelbegin \DIFdel{cannot }\DIFdelend \DIFaddbegin \DIFadd{does not }\DIFaddend distinguish between incomplete aggregation of the smallest particles and limitations of the imaging and tracking workflow at small sizes. \DIFdelbegin \DIFdel{The inferred $d_p$ therefore represents an effective primary scale associated with the population of flocs contributing to the measured settling velocities, rather than the full input grain-size distribution. Because $d_p$ is obtained from regression-line intersections, an apparent truncation of the fine tail reflects a concentration of the mass contributing to the observed $w_s$--$d_f$ scaling within a restricted size range under the experimental freshwater chemistry and shear conditions. The offset between input and inferred distributions may }\DIFdelend \DIFaddbegin \DIFadd{Apparent offsets may also }\DIFaddend arise from size-dependent participation in the floc population resolved by the measurements and from under-sampling of dispersed particles below the detection threshold.

The \DIFaddbegin \DIFadd{absolute magnitude of the inferred $d_p$ is sensitive to the assumed value of the }\DIFaddend shape factor $b_1$\DIFdelbegin \DIFdel{may vary with fractal structure, which would cause $d_\text{intersect}$ to deviate }\DIFdelend \DIFaddbegin \DIFadd{. Laboratory and field studies indicate that $b_1$ is not constant and may increase with $n_f$ }\cite{strom2011explicit}\DIFadd{. If $b_1$ varies from $\mathcal{O}(10)$ to $\mathcal{O}(10^2)$ across fractal classes, the inferred intersection diameter $d_{\mathrm{intersect}}$ can differ }\DIFaddend from the true $d_p$ \DIFdelbegin \DIFdel{. }\DIFdelend \DIFaddbegin \DIFadd{by up to an order of magnitude for fixed $n_f$ (Eq.~\eqref{eq:dpintersection}). However, because the same $b_1(n_f)$ relationship applies uniformly across all experiments, the relative trends in $d_p$ with shear velocity and organic matter reported above are preserved regardless of the assumed shape-factor model; only the absolute scale shifts. 
}

\begin{figure}[htbp]
    \centering
    \DIFaddFL{\includegraphics[width=1\textwidth]{figures/fig11.pdf}
    }\caption{
    \DIFaddFL{Influence of the shape factor (\(b_1\)) on the geometric intersection diameter (\(d_\text{intersect}\)) between floc groups of different fractal dimension.
    (a) Increasing \(b_1\) with \(n_f\) shifts \(d_\text{intersect}\) above the primary particle diameter \(d_p\), while decreasing \(b_1\) shifts it below.
    (b) When evaluating the idealized intersection assuming constant \(b_1\), the lower tail of measured floc diameters extends to the left of \(d_\text{intersect}\). This violates the physical constraint that flocs cannot be smaller than their fundamental building blocks.
    (c) Example where \(d_\text{intersect} > d_p\), emphasizing that \(b_1\) must rise with \(n_f\) to preserve the proper size hierarchy as denser flocs experience greater drag.
    }}
    \label{fig:11}
\end{figure}

\DIFaddend To illustrate this effect, we plot settling velocity versus \DIFaddbegin \DIFadd{floc }\DIFaddend diameter for two \DIFaddbegin \DIFadd{hypothetical }\DIFaddend fractal groups ($n_f = 1.5$ and $n_f = 2.0$), both with a common primary particle diameter of $10~\mu$m, using Eq.~\eqref{eq:floc_settling} (Fig.~\DIFdelbegin \DIFdel{\ref{fig:10}}\DIFdelend \DIFaddbegin \DIFadd{\ref{fig:11}}\DIFaddend a). For $n_f = 1.5$, we fix $b_1 = 50$, and for $n_f = 2.0$, we vary $b_1$ from 10 to 100. As the shape factor increases with fractal dimension, the intersection diameter \DIFdelbegin \DIFdel{$d_\text{intersect}$ }\DIFdelend \DIFaddbegin \DIFadd{$d_{\mathrm{intersect}}$ }\DIFaddend between the two groups becomes greater than the true primary particle diameter. The opposite holds if $b_1$ decreases with $n_f$. 

\DIFdelbegin %DIFDELCMD < \begin{figure}[htbp]
%DIFDELCMD <     \centering
%DIFDELCMD <     %%%
\DIFdelFL{\includegraphics[width=1\textwidth]{figures/10.pdf}
    }%DIFDELCMD < \caption{%
{%DIFAUXCMD
\DIFdelFL{Influence of the shape factor (\(b_1\)) on the geometric intersection diameter (\(d_\text{intersect}\)) between floc groups of different fractal dimension.
    (a) Increasing \(b_1\) with \(n_f\) shifts \(d_\text{intersect}\) above the primary particle diameter \(d_p\), while decreasing \(b_1\) shifts it below.
    (b) When evaluating the idealized intersection assuming constant \(b_1\), the lower tail of measured floc diameters extends to the left of \(d_\text{intersect}\). This violates the physical constraint that flocs cannot be smaller than their fundamental building blocks.
    (c) Example where \(d_\text{intersect} > d_p\), emphasizing that \(b_1\) must rise with \(n_f\) to preserve the proper size hierarchy as denser flocs experience greater drag.
    }}
    %DIFAUXCMD
%DIFDELCMD < \label{fig:10}
%DIFDELCMD < \end{figure}
%DIFDELCMD < 

%DIFDELCMD < %%%
\DIFdelend We evaluate the calculated geometric intersection diameters ($d_\text{intersect}$) obtained under the standard assumption of a constant shape factor (Fig.~\DIFdelbegin \DIFdel{\ref{fig:10}}\DIFdelend \DIFaddbegin \DIFadd{\ref{fig:11}}\DIFaddend b). \DIFdelbegin \DIFdel{We find that while }\DIFdelend \DIFaddbegin \DIFadd{While }\DIFaddend most measured floc diameters lie to the right of $d_\text{intersect}$, the lower tail of the observed floc size distribution inevitably extends to the left (i.e., some observed $d_f < d_\text{intersect}$). Because a floc physically cannot be smaller than its constituent primary particles ($d_f \ge d_p$), the true primary particle diameter must be strictly smaller than all measured floc diameters. Consequently, the assumption that $d_\text{intersect} = d_p$ produces a physical contradiction, requiring instead that the true primary particle size be securely anchored to the left of the constant-$b_1$ intersection ($d_p < d_\text{intersect}$). 

This strict physical constraint implies a necessary relationship between $b_1$ and $n_f$. According to Eq.~\eqref{eq:dpintersection}, the only way to \DIFdelbegin \DIFdel{mathematically }\DIFdelend maintain $d_p < d_\text{intersect}$ as $n_f$ increases is if $b_1$ also increases. \DIFdelbegin \DIFdel{This requirement ensures that the primary particle remains the smallest logical constituent in the system. }\DIFdelend Such a trend is consistent with theoretical expectations and experimental observations, which suggest that denser, less permeable flocs (i.e., those with higher $n_f$) experience greater hydrodynamic drag, corresponding to higher values of $b_1$ \cite{strom2011explicit}. Importantly, this conclusion does not invalidate the $n_f$ values inferred from the stratified $\log(w_s)$--$\log(d_f)$ slopes, because those slopes are independent of $b_1$; the shape factor enters only the intercept. The sensitivity of $d_\text{intersect}$ to $b_1$ therefore affects only the \DIFaddbegin \DIFadd{magnitude of the }\DIFaddend inferred primary particle diameter, not the fractal dimension itself \DIFaddbegin \DIFadd{nor the systematic variation of $d_p$ across experimental conditions}\DIFaddend .

\DIFdelbegin \section{\DIFdel{Discussion}}
%DIFAUXCMD
\addtocounter{section}{-1}%DIFAUXCMD
%DIFDELCMD < 

%DIFDELCMD < %%%
\DIFdelend \subsection{\DIFdelbegin \DIFdel{How Fractal Structure Mediates }\DIFdelend \DIFaddbegin \DIFadd{Reconciling }\DIFaddend Settling \DIFdelbegin \DIFdel{Velocity}\DIFdelend \DIFaddbegin \DIFadd{Dynamics with Variable Fractal Dimension}\DIFaddend }
\DIFaddbegin \label{sec:reconciling}
\DIFaddend 

\DIFdelbegin \DIFdel{A central finding of this study is that the 3D fractal dimension varies systematically with environmental forcing and that this variation has direct consequences for settling velocity . The stratified analysis reveals that }\DIFdelend \DIFaddbegin \DIFadd{Having established that both }\DIFaddend $n_f$ \DIFdelbegin \DIFdel{increases with shear velocity (from 1.77 to 2.12 over }\DIFdelend \DIFaddbegin \DIFadd{(Section~\ref{sec:stratifying}) and $d_p$ (Section~\ref{sec:dp}) vary systematically with experimental conditions, we now test whether this structural variability explains the settling velocity discrepancies identified in Section~\ref{sec:bulk}. We formulate a modified model that retains the same equilibrium $d_f$ prediction but evaluates $w_s \propto d_p^{\,3-n_f}\, d_f^{\,n_f-1}$ with the measured, condition-specific $n_f$ and $d_p$ rather than fixed values. As with the base model, we normalize the prediction by a single least-squares constant and plot it against the measured medians in Fig.~\ref{fig:6}c,d. However, both models share the identical $d_{f,\text{eq}}$ prediction and each contain exactly one free parameter, so any improvement in the modified model arises solely from incorporating the measured structural variability, not from additional fitting degrees of freedom.
}

\DIFadd{As shown in Fig.~\ref{fig:6}c, the modified model successfully reproduces the non-monotonic peak in settling velocity with shear. At low to moderate shear (}\DIFaddend $u_* = 0.02$ \DIFdelbegin \DIFdel{--$0.04$}\DIFdelend \DIFaddbegin \DIFadd{to $0.03$~m\,s$^{-1}$), the rapid densification of flocs (increasing $n_f$) outpaces the reduction in aggregate size, causing $w_s$ to rise despite shrinking $d_f$. Only at high shear ($u_* = 0.04$}\DIFaddend ~m\,s$^{-1}$) \DIFdelbegin \DIFdel{and decreases }\DIFdelend \DIFaddbegin \DIFadd{does fragmentation dominate over densification, and $w_s$ declines. Similarly, the modified model captures the steep monotonic decrease of $w_s$ }\DIFaddend with organic matter \DIFdelbegin \DIFdel{content (from 1.84 to 1.63 over 1--3~wt\% OM). These trends reflect two competing physical processes that govern floc structure: shear-driven breakup, which preferentially destroys fragile, open aggregates and selects for compact survivors with higher $n_f$; and organic binding, in which extracellular polymeric substances act as a flexible ``glue'' that holds particles in extended, porous configurations with lower }\DIFdelend \DIFaddbegin \DIFadd{(Fig.~\ref{fig:6}d): decreasing }\DIFaddend $n_f$ \DIFdelbegin %DIFDELCMD < \cite{nghiem2022mechanistic}%%%
\DIFdelend \DIFaddbegin \DIFadd{and $d_p$ with added OM reduce the effective solid fraction of the aggregates, overwhelming the minor increase in $d_f$}\DIFaddend .

\DIFdelbegin \DIFdel{These structural changes propagate directly into settling velocity }\DIFdelend \DIFaddbegin \DIFadd{Across all six experiments, the modified model reduces the root-mean-square error in predicted settling velocity by 34\% relative to the base model. This improvement demonstrates that the discrepancies between the constant-$n_f$ framework and our data are not random scatter but arise from systematic structural changes that the base model cannot represent. Properly predicting the settling dynamics of cohesive sediment therefore requires accounting for the dynamic variability of $n_f$ and $d_p$ with environmental forcing.
}

\section{\DIFadd{Discussion}}

\DIFadd{A central finding of this research is the discrepancy between fractal dimensions derived from aggregated (i.e., pooled) datasets and those obtained from datasets stratified by 2D box-counting dimension. As illustrated by our conceptual model (Fig.~\ref{fig:1}) and supported by experimental results (Figs.~\ref{fig:8} and \ref{fig:9}), aggregating settling data from a floc population with inherent structural variability---meaning a range of true $n_f$ values---can yield a bulk-fitted $n_f$ that overrepresents certain structural subgroups while failing to capture the full diversity within the population. This approach can obscure underlying trends and produce fitted fractal-dimension values that are potentially misleading. The mechanism is analogous to Simpson’s paradox: because distinct structural subpopulations occupy different regions of $w_s$--$d_f$ space, the slope of a pooled regression is not a simple average of within-group slopes and may be dominated by the subpopulation that exerts the greatest leverage in log--log space}\DIFaddend .
\DIFdelbegin \DIFdel{From }\DIFdelend \DIFaddbegin 

\subsection{\DIFadd{Effect of Shear Velocity and Organic Matter}}
\DIFadd{The settling velocity $w_s$ governs the rate at which flocculated mud deposits from the water column and is therefore the quantity of most direct relevance to sediment transport predictions. For fractal flocs, $w_s \propto d_p^{\,3-n_f}\, d_f^{\,n_f-1}$ (}\DIFaddend Eq.~\eqref{eq:floc_settling}\DIFdelbegin \DIFdel{, $w_s$ depends on floc size as $d_f^{n_f-1}$ and on floc density (via $d_p^{3-n_f}$), so both sizeand $n_f$ jointly determine the settling rate.
In the shear series, increasing $u_*$ simultaneously decreases floc size (median }\DIFdelend \DIFaddbegin \DIFadd{), which can be rewritten as $w_s \propto \phi \, d_f^2$ by defining the solid volume fraction $\phi = (d_f / d_p)^{n_f - 3}$. In this form, settling is controlled jointly by floc size ($d_f^2$) and effective density ($\phi$), where the latter encapsulates the combined influence of fractal dimension and primary particle size. This decomposition is useful because shear velocity and organic matter affect }\DIFaddend $d_f$ \DIFdelbegin \DIFdel{drops from 43 to 27~}%DIFDELCMD < \textmu %%%
\DIFdel{m; Table~\ref{tab:bulk_trends})and increases compactness ($n_f$ rises from 1.77 to 2.12). These effects partially offset each other: smaller flocs settle more slowly, but denser flocs settle faster.The net result is that median $w_s$ remains roughly comparable across shear levels ($0.30$--$0.39$~mm\,s$^{-1}$}\DIFdelend \DIFaddbegin \DIFadd{and $\phi$ in distinct and sometimes opposing ways.
}

\DIFadd{With increasing shear velocity, $d_f$ decreases as stronger turbulence breaks apart aggregates (Fig.~\ref{fig:6}a), while $n_f$ increases (Fig.~\ref{fig:9}a) and the inferred $d_p$ coarsens (Fig.~\ref{fig:10}d}\DIFaddend ), consistent with \DIFdelbegin \DIFdel{the non-monotonic bulk trend in Table~\ref{tab:bulk_trends}. In the OM series, increasing organic content both reduces compactness (}\DIFdelend \DIFaddbegin \DIFadd{preferential retention of compact, shear-resistant structures }\cite{winterwerp1998simple}\DIFadd{. The rising }\DIFaddend $n_f$ \DIFdelbegin \DIFdel{drops from 1.84 to 1.63) and modestly decreases floc size (median }\DIFdelend \DIFaddbegin \DIFadd{and coarsening $d_p$ increase $\phi$, partially compensating for the shrinking }\DIFaddend $d_f$\DIFdelbegin \DIFdel{drops from 42 to 39~}%DIFDELCMD < \textmu %%%
\DIFdel{m)}\DIFdelend \DIFaddbegin \DIFadd{. Because these two contributions act in opposing directions on $w_s$, the observed settling response is non-monotonic (Fig.~\ref{fig:6}c). With increasing organic matter, $n_f$ decreases (Fig.~\ref{fig:9}b) and $d_p$ fines (Fig.~\ref{fig:10}h), consistent with organic binding promoting open, porous architectures built from smaller primary particles}\DIFaddend . Both effects reduce \DIFdelbegin \DIFdel{$w_s$, and their reinforcement drives the pronounced decline from 0.78 to 0.52~mm\,s$^{-1}$}\DIFdelend \DIFaddbegin \DIFadd{$\phi$, and because $d_f$ remains relatively stable across the OM gradient (Fig.~\ref{fig:6}b), $w_s$ decreases monotonically (Fig.~\ref{fig:6}d)}\DIFaddend .

\DIFdelbegin \DIFdel{In contrast, the aggregated (bulk) approach yields $n_f \approx 1.9$--$2.0$ }\DIFdelend \DIFaddbegin \DIFadd{The equilibrium floc diameter $d_{f,\text{eq}}$, by contrast, is well predicted by the }\citeA{nghiem2022mechanistic} \DIFadd{framework }\DIFaddend across all conditions \DIFdelbegin \DIFdel{and would predict that changes in $w_s$ are driven almost entirely by size.The stratified analysis reveals that this interpretation is incomplete: structural changes account for asubstantial fraction of the settling-velocity response,particularly in the OM series where the stratified }\DIFdelend \DIFaddbegin \DIFadd{(Fig.~\ref{fig:6}a,b), because }\DIFaddend $n_f$ \DIFdelbegin \DIFdel{drops well below the aggregated estimate.
}%DIFDELCMD < 

%DIFDELCMD < %%%
\DIFdel{These trends are broadly consistent with prior work. }%DIFDELCMD < \citeA{winterwerp1998simple} %%%
\DIFdel{reported }\DIFdelend \DIFaddbegin \DIFadd{cancels in the aggregation--breakup balance and the predicted size depends only on $\eta$ and $\theta$. The settling velocity retains an explicit dependence on }\DIFaddend $n_f$ \DIFdelbegin \DIFdel{as low as 1.4 for fragile marine-snow aggregates and up to 2.2 for robust estuarine flocs, and attributed the range to differences in shear history and biological binding.Our results extend this framework by quantifying, within a single experimental system,the rates at which }\DIFdelend \DIFaddbegin \DIFadd{and $d_p$, and a constant-}\DIFaddend $n_f$ \DIFdelbegin \DIFdel{responds to shear ($\approx$17~s\,m$^{-1}$) and organic matter ($\approx -0.11$ per wt\% ), and by demonstrating that these responses are detectable only when structural heterogeneity is resolved}\DIFdelend \DIFaddbegin \DIFadd{model cannot reproduce either the non-monotonic shear response or the organic-matter decline (Fig.~\ref{fig:6}c,d). Incorporating the measured structural variability (Section~\ref{sec:reconciling}) reduces the root-mean-square error by 34\% with no additional free parameters, confirming that these discrepancies reflect systematic structural changes rather than scatter. Although the exact quantitative relationships are likely system-specific, the general trends are broadly consistent with previous studies and underscore that accurate settling predictions require accounting for the variability of $n_f$ and $d_p$ with environmental forcing}\DIFaddend .

\subsection{2D to 3D Conversion of Fractal Dimension}
\DIFdelbegin %DIFDELCMD < 

%DIFDELCMD < %%%
\DIFdel{The stratified data reveal a positive, nonlinear relationship between }\DIFdelend \DIFaddbegin \DIFadd{The relationship between the }\DIFaddend 2D box-counting \DIFdelbegin \DIFdel{fractal dimension and the hydrodynamically }\DIFdelend \DIFaddbegin \DIFadd{dimension ($n_f^{(\mathrm{2D})}$) and the }\DIFaddend inferred 3D fractal dimension (\DIFaddbegin \DIFadd{$n_f$) is broadly consistent with models developed for synthetic computational aggregates (}\DIFaddend Fig.~\DIFdelbegin \DIFdel{\ref{fig:8}c,d}\DIFdelend \DIFaddbegin \DIFadd{\ref{fig:9}c}\DIFaddend ). The \DIFdelbegin \DIFdel{idealized }\DIFdelend \DIFaddbegin \DIFadd{simple }\DIFaddend Euclidean conversion, \DIFdelbegin \DIFdel{$n_f = n_f^{(2\mathrm{D})} + 1$, systematically overpredicts }\DIFdelend \DIFaddbegin \DIFadd{$n_f = n_f^{(\mathrm{2D})} + 1$, systematically overestimates }\DIFaddend $n_f$ \DIFdelbegin \DIFdel{because porous aggregates do not fill space completely: a 2D projection of a fractal object loses structural information, so the 3D dimension grows sublinearly. Our data align most }\DIFdelend \DIFaddbegin \DIFadd{across our measured range, whereas our data align more }\DIFaddend closely with the empirical \DIFdelbegin \DIFdel{box-counting conversion }\DIFdelend \DIFaddbegin \DIFadd{conversion curves }\DIFaddend of \citeA{wang2022fractal}\DIFdelbegin \DIFdel{($n_f = 0.81\,(n_f^{(2\mathrm{D})})^{1.81}$), derived from }\DIFdelend \DIFaddbegin \DIFadd{, based on }\DIFaddend computationally generated aggregates. 

\DIFdelbegin \DIFdel{Although this agreement is encouraging}\DIFdelend \DIFaddbegin \DIFadd{However}\DIFaddend , no universally valid mapping between 2D and 3D fractal dimensions exists for natural flocs. \DIFdelbegin \DIFdel{Synthetic }\DIFdelend \DIFaddbegin \DIFadd{Most synthetic }\DIFaddend studies, including \citeA{wang2022fractal}, grow aggregates from monodisperse primary particles with fixed size, \DIFdelbegin \DIFdel{whereas natural flocs}\DIFdelend \DIFaddbegin \DIFadd{allowing simple scaling laws to be derived. Natural flocs, by contrast, }\DIFaddend contain a broad distribution of particle sizes\DIFdelbegin \DIFdel{and shapes, violating the fixed-$d_p$ assumption}\DIFdelend \DIFaddbegin \DIFadd{, shapes, and compositions, adding heterogeneity in packing and porosity that is not captured by idealized models. Other approaches }\cite{maggi2004method, tang2015reconstructing} \DIFadd{depend on their own specific assumptions and likewise cannot account for this variability}\DIFaddend . These differences introduce scatter and \DIFdelbegin \DIFdel{potential systematic offsets . Furthermore, permeability---which is not considered here---can alter settling without changing geometry, adding an additional unmapped degree of freedom. The 2D--3D relationship should therefore be treated as system specific and re-evaluated when sediment mineralogy, salinity, or aggregation pathways differ substantially}\DIFdelend \DIFaddbegin \DIFadd{systematic offsets between natural data and synthetic conversions.
}

\DIFadd{The close match between our results and the box-counting conversion of }\citeA{wang2022fractal} \DIFadd{may therefore reflect structural similarities between our experimental flocs and their synthetic aggregates rather than a universally applicable rule}\DIFaddend . Nevertheless, the 2D box-counting dimension remains a practical stratification tool: by quantifying within-population variability \DIFdelbegin \DIFdel{and enabling direct experimental calibration}\DIFdelend \DIFaddbegin \DIFadd{in $n_f^{(\mathrm{2D})}$ and calibrating the 2D--3D relationship directly from experimental data}\DIFaddend , our approach provides context-sensitive inference of the hydrodynamically relevant 3D fractal dimension \DIFaddbegin \DIFadd{without relying on a fixed conversion law}\DIFaddend .

\subsection{Constraining Primary Particle Diameter ($d_p$) and Shape Factor ($b_1$)}
\DIFdelbegin %DIFDELCMD < 

%DIFDELCMD < %%%
\DIFdel{The primary particle diameter ($d_p$) and the shape factor ($b_1$) are key parameters in fractal settling theory, yet they are rarely measured directly. We estimate $d_p$ by identifying $d_\text{intersect}$, the diameter at which the $n_f$-stratified log--log regression lines converge (Fig.~\ref{fig:9}). In principle, at $d_f = d_p$ all floc populations should converge to the settling velocity of a single primary particle regardless of $n_f$.
}%DIFDELCMD < 

%DIFDELCMD < %%%
\DIFdelend The intersection analysis yields \DIFdelbegin \DIFdel{median }\DIFdelend $d_\text{intersect}$ \DIFdelbegin \DIFdel{values that agree with the median of the input sediment sizes }\DIFdelend \DIFaddbegin \DIFadd{distributions that agree in central tendency with the input sediment but are substantially narrower, peaking in the coarse-silt range }\DIFaddend (Fig.~\DIFdelbegin \DIFdel{\ref{fig:9}), yet the inferred distribution is narrower: about two orders of magnitude wide compared with nearly four orders in the input material. This discrepancy may arise because only a subset of the source material (coarse silt to fine sand) effectively participates in forming the flocs that dominate the measured settling velocities, or because the imaging system does not capture the smallest dispersed grains}\DIFdelend \DIFaddbegin \DIFadd{\ref{fig:10}). This narrowing is expected on theoretical grounds}\DIFaddend . \citeA{nghiem2024testing} showed that the effective $d_p$ in fractal scaling \DIFaddbegin \DIFadd{theory }\DIFaddend is better represented by a number-weighted moment of the size distribution rather than a volumetric median, because \DIFdelbegin \DIFdel{physical }\DIFdelend \DIFaddbegin \DIFadd{fractal }\DIFaddend aggregation conserves particle number rather than volume. The inferred \DIFaddbegin \DIFadd{$d_\text{intersect}$ is therefore expected to deviate from the bulk median grain size in a manner that depends on the weighting imposed by the aggregation process. Additionally, the finest clays may remain dispersed under the experimental freshwater chemistry, while the coarsest silica grains may be too heavy to be incorporated into flocs---both effects narrow the effective }\DIFaddend $d_p$ \DIFdelbegin \DIFdel{values therefore represent an effective primary scale for the resolved floc population rather than a census of all grains present}\DIFdelend \DIFaddbegin \DIFadd{distribution relative to the input}\DIFaddend . 

\DIFdelbegin \DIFdel{A key result of this analysis is the demonstration that }\DIFdelend \DIFaddbegin \DIFadd{The intersection method assumes a constant }\DIFaddend $b_1$ \DIFdelbegin \DIFdel{must increase with }\DIFdelend \DIFaddbegin \DIFadd{across }\DIFaddend $n_f$ \DIFdelbegin \DIFdel{. Under the assumption of constant }\DIFdelend \DIFaddbegin \DIFadd{classes (Eq.~\eqref{eq:dpintersectionideal}); when }\DIFaddend $b_1$ \DIFdelbegin \DIFdel{, the intersection diameter }\DIFdelend \DIFaddbegin \DIFadd{varies with floc structure, }\DIFaddend $d_\text{intersect}$ \DIFdelbegin \DIFdel{exceeds some observed floc diameters (Fig.~\ref{fig:10}b), which is physically impossible because flocs cannot be smaller than their building blocks. }\DIFdelend \DIFaddbegin \DIFadd{deviates from the true $d_p$ (}\DIFaddend Eq.~\DIFdelbegin \DIFdel{\eqref{eq:dpintersection} shows that $d_\text{intersect} > d_p$ requires }\DIFdelend \DIFaddbegin \DIFadd{\eqref{eq:dpintersection}). Our analysis (Fig.~\ref{fig:11}) shows that physically plausible intersections require }\DIFaddend $b_1$ to increase with $n_f$\DIFdelbegin \DIFdel{: }\DIFdelend \DIFaddbegin \DIFadd{, consistent with the expectation that }\DIFaddend denser, less permeable flocs experience greater form drag \DIFdelbegin \DIFdel{and thus higher effective $b_1$ }\DIFdelend \cite{strom2011explicit}. \DIFdelbegin \DIFdel{This finding is consistent with theoretical expectations but, to our knowledge, has not previously been demonstrated empirically from settling data alone. We emphasize that this $b_1$--$n_f$ dependence }\DIFdelend \DIFaddbegin \DIFadd{However, this sensitivity }\DIFaddend affects only the \DIFdelbegin \DIFdel{inferred $d_p$ (via the regression intercept) and does not alter the slope-based }\DIFdelend \DIFaddbegin \DIFadd{absolute magnitude of $d_\text{intersect}$. The }\DIFaddend $n_f$ \DIFdelbegin \DIFdel{estimates, which }\DIFdelend \DIFaddbegin \DIFadd{values inferred from stratified slopes }\DIFaddend are independent of $b_1$\DIFaddbegin \DIFadd{, which enters only the intercept, and the relative trends in $d_\text{intersect}$ across experimental conditions are preserved because the same $b_1(n_f)$ relationship applies uniformly to all experiments. Independently constraining $b_1(n_f)$---for example through direct drag measurements on size-selected flocs---would allow the intersection estimate to be converted into a true $d_p$}\DIFaddend . 

\DIFdelbegin \DIFdel{Uncertainty in }\DIFdelend \DIFaddbegin \DIFadd{Floc permeability, not considered here, may further complicate the inference. Permeability reduces the effective drag coefficient and may itself depend on $d_f$ and }\DIFaddend $d_p$\DIFdelbegin \DIFdel{propagates directly into predictions: the value of }\DIFdelend \DIFaddbegin \DIFadd{; if it varies systematically with floc size, assuming it to be constant can bias both the inferred slope and therefore }\DIFaddend $n_f$ \DIFdelbegin \DIFdel{inferred from a full-model fit to Eq. ~\eqref{eq:floc_settling} is highly sensitive to the assumed $d_p$ (Fig.~\ref{fig:4}c). 
The slope-based approach adopted hereavoids this coupling, but predictive use of Eq.~\eqref{eq:floc_settling} ultimately requires both parameters. In practice, the effective $d_p$ relevant for flocculation can be both narrower in range and offset in value relative to the bulk sediment size distribution}\DIFdelend \DIFaddbegin \cite{nghiem2024testing}\DIFadd{. Quantifying the joint dependence of $b_1$ and permeability on floc structure remains an important target for future work}\DIFaddend .

\subsection{Implications\DIFdelbegin \DIFdel{for Sediment Transport Modeling}\DIFdelend }
\DIFdelbegin %DIFDELCMD < 

%DIFDELCMD < %%%
\DIFdelend The 3D fractal dimension \DIFaddbegin \DIFadd{of mud flocs, }\DIFaddend $n_f$\DIFdelbegin \DIFdel{governs the density--size relationship of flocs and thereby controls settling velocity }\DIFdelend \DIFaddbegin \DIFadd{, is a key parameter }\DIFaddend in morphodynamic models \DIFaddbegin \DIFadd{of fine-grained sediment transport}\DIFaddend . Historically, modelers have \DIFdelbegin \DIFdel{adopted $n_f \approx 2.0$--$2.2$ from bulk $w_s$--$d_f$ regressions }%DIFDELCMD < \cite{kranenburg1994fractal, son2008flocculation, winterwerp1998simple, winterwerp2004introduction, winterwerp2006heuristic}%%%
\DIFdel{. Our aggregated analysis yields $n_f \approx 2.0$, confirming }\DIFdelend \DIFaddbegin \DIFadd{generally assumed $n_f \approx 2$ }\cite{kranenburg1994fractal, son2008flocculation, winterwerp1998simple, winterwerp2004introduction}\DIFadd{. These studies inferred the fractal dimension by aggregating settling velocity data from entire floc populations. When we apply this same classical aggregation technique to our own experimental data, we derive a similar fractal dimension of approximately 2.0. This consistency confirms }\DIFaddend that our experimental flocs are physically comparable to those \DIFdelbegin \DIFdel{in previous studies}\DIFdelend \DIFaddbegin \DIFadd{observed in previous work}\DIFaddend . However, \DIFdelbegin \DIFdel{the stratified analysis reveals that this apparent consistency masks substantial structural diversity}\DIFdelend \DIFaddbegin \DIFadd{this apparent agreement masks a significant underlying complexity}\DIFaddend . 

By \DIFdelbegin \DIFdel{resolving this diversity, the stratified method yields }\DIFdelend \DIFaddbegin \DIFadd{stratifying flocs according to their 2D box-counting dimension, our method reveals systematic trends that are invisible to conventional data aggregation. We find that the }\DIFaddend median $n_f$ \DIFdelbegin \DIFdel{values of 1.6--2.1 depending on forcing, with the majority of conditions producing $n_f < 2.0$. These lower fractal dimensions correspond to more porous, less dense flocs that settle more slowly. For a representative }\DIFdelend \DIFaddbegin \DIFadd{increases with shear velocity (from 1.77 to 2.12) and decreases with increasing organic matter content (from 1.84 to 1.63). When the stratified $n_f$ and $d_p$ values are fed back into the settling velocity equation (Section~\ref{sec:reconciling}), the modified model reduces the prediction error by 34\% relative to the constant-$n_f$ base model, providing direct evidence that this structural variability is dynamically significant. In contrast, the aggregated analysis compresses the dynamic range by roughly a factor of five, yielding median values that vary by only $\sim$0.06 across the same experimental conditions compared with $\sim$0.2--0.35 for the stratified approach. The aggregated estimates also regress toward a population-average $n_f \approx 2$, reflecting Simpson’s paradox: fitting a single slope to a structurally mixed population of loose and dense flocs yields an intermediate value that neither represents the dominant subgroup nor captures the underlying environmental trend. The result is that relying on a single, representative fractal dimension near 2.0, as is typical in existing models, both overstates floc compactness under conditions that produce porous flocs and obscures the sensitivity of $n_f$ to shear and organic matter. 
}

\DIFadd{This finding has important implications for sediment transport prediction. Using the Stokes' law--based formulation for floc settling velocity, $w_s \propto d_p^{3-n_f}\, d_f^{n_f-1}$, lowering $n_f$ from 2.0 to 1.7 for a typical }\DIFaddend 400~\textmu m floc \DIFaddbegin \DIFadd{(}\DIFaddend with 4~\textmu m primary particles, \DIFdelbegin \DIFdel{lowering $n_f$ from 2.2 to 1.7 in Eq.~\eqref{eq:floc_settling} produces an order-of-magnitude }\DIFdelend \DIFaddbegin \DIFadd{a representative value from the literature) produces a several-fold }\DIFaddend decrease in settling \DIFdelbegin \DIFdel{velocity, and a corresponding order-of-magnitude }\DIFdelend \DIFaddbegin \DIFadd{speed, which directly translates to a several-fold }\DIFaddend increase in transport distance. In many transport formulations \DIFdelbegin \DIFdel{, }\DIFdelend the deposition flux scales as \DIFdelbegin \DIFdel{$F_{\mathrm{dep}} \sim w_s\,C$ (with suspended concentration $C$) }\DIFdelend \DIFaddbegin \DIFadd{$F_\text{dep} \sim w_s\, C$ (where $C$ is the suspended sediment concentration) }\cite{amoudry2008review}\DIFaddend , so lowering $n_f$\DIFdelbegin \DIFdel{and }\DIFdelend \DIFaddbegin \DIFadd{---and }\DIFaddend thus $w_s$\DIFdelbegin \DIFdel{reduces }\DIFdelend \DIFaddbegin \DIFadd{---reduces }\DIFaddend modeled deposition rates and increases downstream export for the same supply\DIFdelbegin %DIFDELCMD < \cite{amoudry2008review}%%%
\DIFdelend \DIFaddbegin \DIFadd{. Flocs that might be predicted to settle out over a given distance would travel several times farther under this revised, more porous structural model. As a result, mud flocs in natural environments may settle more slowly and be transported substantially farther than predicted by existing models that rely on aggregated fractal dimensions}\DIFaddend . The efficiency of mud retention\DIFdelbegin \DIFdel{in freshwater floodplains and deltas may therefore }\DIFdelend \DIFaddbegin \DIFadd{, at least in the freshwater environments considered here, may }\DIFaddend be considerably lower than \DIFdelbegin \DIFdel{models assuming $n_f \geq 2.0$ would predict.
}%DIFDELCMD < 

%DIFDELCMD < %%%
\DIFdel{For practitioners, we recommend adopting $n_f$ in the range of 1.6--2.0, and, where feasible, representing the fractal dimension as a distribution rather than a single value. Selecting a lower $n_f$ yields predictions of reduced deposition and increased sediment bypass relative to the conventional range. While next-generation models should ideally represent the full distribution of $n_f$ to capture natural heterogeneity, a pragmatic immediate step is to shift default parameterizations away from the biased, artificially high values that arise from bulk data aggregation. This adjustment is important for improving predictions of fine-sediment fate in freshwater systems; extending the framework to saline estuaries requires explicit consideration of salt-driven physical aggregation, which may shift the $n_f$ distribution toward higher values}\DIFdelend \DIFaddbegin \DIFadd{currently assumed, with significant consequences for sediment budgets, contaminant dispersal, and carbon burial in floodplains and deltas}\DIFaddend .

\section{\DIFdelbegin \DIFdel{Conclusions}\DIFdelend \DIFaddbegin \DIFadd{Conclusion}\DIFaddend }
\DIFdelbegin %DIFDELCMD < 

%DIFDELCMD < %%%
\DIFdelend This study demonstrates that accounting for \DIFdelbegin \DIFdel{structural }\DIFdelend heterogeneity within floc populations is essential for accurate sediment transport modeling. Aggregating settling data across flocs with varying fractal dimensions \DIFdelbegin \DIFdel{produces }\DIFdelend \DIFaddbegin \DIFadd{can produce }\DIFaddend misleading results due to Simpson's \DIFdelbegin \DIFdel{Paradox: physical }\DIFdelend \DIFaddbegin \DIFadd{paradox: }\DIFaddend trends present within \DIFdelbegin \DIFdel{distinct }\DIFdelend structural subgroups are lost or distorted when data are \DIFdelbegin \DIFdel{analyzed in bulk}\DIFdelend \DIFaddbegin \DIFadd{pooled}\DIFaddend . By stratifying flocs based on \DIFaddbegin \DIFadd{their }\DIFaddend 2D box-counting fractal \DIFdelbegin \DIFdel{dimensions}\DIFdelend \DIFaddbegin \DIFadd{dimension}\DIFaddend , we resolve \DIFdelbegin \DIFdel{these structural differences and enable physically meaningful inference of floc transport parameters}\DIFdelend \DIFaddbegin \DIFadd{structural differences that conventional bulk analysis obscures}\DIFaddend . 

The stratified approach reveals that \DIFdelbegin \DIFdel{the hydrodynamically inferred 3D fractal dimension (}\DIFdelend $n_f$ \DIFdelbegin \DIFdel{) }\DIFdelend increases with shear velocity \DIFdelbegin \DIFdel{($u_*$) }\DIFdelend and decreases with organic matter content, \DIFdelbegin \DIFdel{reflecting the competition between shear-driven breakup---which selects for compact, high-}\DIFdelend \DIFaddbegin \DIFadd{while the inferred $d_p$ co-varies with these same conditions---coarsening under stronger shear and fining under organic binding. The concurrent variation of }\DIFaddend $n_f$\DIFdelbegin \DIFdel{structures---and organic binding, which promotes open, porous aggregates with low $n_f$. These trends, quantified as $\Delta n_f / \Delta u_* \approx 17$~s\,m$^{-1}$ and $\Delta n_f / \Delta \mathrm{OM} \approx -0.11$ per wt\%, are invisible in bulk analyses that yield a near-constant $n_f \approx 2.0$. The relationship between 2D }\DIFdelend \DIFaddbegin \DIFadd{, $d_p$, }\DIFaddend and \DIFdelbegin \DIFdel{3D fractal dimensions is inherently nonlinear, aligning with the empirical box-counting conversion of }%DIFDELCMD < \citeA{wang2022fractal} %%%
\DIFdel{rather than the idealized Euclidean offset, consistent with the sub-Euclidean space filling of porous aggregates. Furthermore, the stratified intersection analysis constrains the effective primary particle diameter and demonstrates that the hydrodynamic shape factor (}\DIFdelend \DIFaddbegin \DIFadd{$d_f$ introduces competing effects on $w_s \propto d_p^{3-n_f}\, d_f^{n_f-1}$, producing non-monotonic settling responses that constant-$n_f$ models cannot capture. Incorporating the measured structural variability into the settling velocity prediction reduces the root-mean-square error by 34\% relative to the constant-$n_f$ base model, with no additional free parameters. The framework also yields an empirical 2D--3D fractal dimension relationship that matches existing synthetic-aggregate conversions and shows that }\DIFaddend $b_1$ \DIFdelbegin \DIFdel{) }\DIFdelend must increase with $n_f$ to maintain physical consistency. 

These results \DIFdelbegin \DIFdel{show that the widespread use of $n_f \approx 2.0$--$2.2$ in sediment transport models }\DIFdelend \DIFaddbegin \DIFadd{indicate that the conventional assumption of $n_f \approx 2$ }\DIFaddend overestimates floc compactness \DIFdelbegin \DIFdel{, leading to overprediction of settling velocity by up to an order of magnitude and corresponding underprediction of transport distance. We recommend adopting }\DIFdelend \DIFaddbegin \DIFadd{and settling rates. Reducing }\DIFaddend $n_f$ \DIFdelbegin \DIFdel{in the range }\DIFdelend \DIFaddbegin \DIFadd{to the observed range of }\DIFaddend 1.6--2\DIFdelbegin \DIFdel{.0, which better represents the porous, slow-settling flocs that dominate suspended transport but remain hidden in bulk analyses. Making this adjustment will improve predictions of fine-sediment transport, deltaic development}\DIFdelend \DIFaddbegin \DIFadd{.1 can decrease predicted $w_s$ by several-fold for typical floc sizes, implying that mud retention in freshwater systems may be considerably lower than current models suggest. We recommend adopting a lower default $n_f$ and, where feasible, representing the fractal dimension as a distribution. While no single value captures the full complexity of natural flocculation, shifting away from the classical range is a necessary step toward more accurate predictions of fine sediment transport, deposition}\DIFaddend , and downstream sediment budgets\DIFdelbegin \DIFdel{in freshwater-dominated systems}\DIFdelend .



\section*{Open Research Section}
All code and analysis workflows supporting the results and conclusions of this study are openly available at \url{https://github.com/braydennoh/FlocTrack}. The repository includes scripts for image analysis, fractal dimension calculation, and all processing steps described in this manuscript. For further questions or access to additional datasets, please contact the corresponding author at braydennoh@fas.harvard.edu.

\section*{Inclusion in Global Research Statement}
This research did not involve collaborations with local communities or cross-regional partnerships requiring additional disclosure. All research activities were conducted in accordance with institutional and international ethical standards.

\acknowledgments
BN acknowledges support from Caltech SURF, the Arthur R. Adams SURF Fellowship, and the Saul and Joan Cogen Memorial SURF Fellowship at Caltech. MPL acknowledges support from NASA FINESST Grant 80NSSC20K1645, the NASA Delta-X project funded by the Science Mission Directorate's Earth Science Division through the Earth Venture Suborbital-3 Program NNH17ZDA001N-EVS3, and the National Science Foundation Geomorphology and Land-use Dynamics Grant No. 2136991.

\bibliography{agusample}


\end{document}
